%%%%%%%%%%%%%%%%%%%%%%%%%%%%%%%%%%%%%%%%%%%%%%%%%%%%%%%%%%%%%%%%%%%%%%
% Lớp 11 | Xác suất | Biến cố giao và quy tắc nhân xác suất
% Dạng 1D9?1-3: Công thức nhân xác suất cho 2 biến cố độc lập
% Số câu: 75  (sinh tự động bởi tach_cau.py; từ nay thầy sửa trực tiếp file này)
%%%%%%%%%%%%%%%%%%%%%%%%%%%%%%%%%%%%%%%%%%%%%%%%%%%%%%%%%%%%%%%%%%%%%%

%%%%%%%%%%%%%%%% TRẮC NGHIỆM 4 PHƯƠNG ÁN %%%%%%%%%%%%%%%%

% nguồn: Ok_TracNghiemLop11_Lan1
\begin{ex}%[1D9N1-3]%[ma:ae32e0a1]
	Cho $A$ và $B$ là hai biến cố độc lập với nhau thỏa mãn $\mathrm{P}(A)=0{,}4$,  $\mathrm{P}(B)=0{,}3$. Khi đó $\mathrm{P}(AB)$ bằng
	 \choice
		{$0{,}58$} 
		{$0{,}7$} 
		 {$0{,}1$}
		{\True $0{,}12$} 
	 \loigiai{
		Do $A$ và  $B$ là hai biến cố độc lập với nhau nên $\mathrm{P}(AB)=\mathrm{P}(A).\mathrm{P}(B)=0{,}4.0{,}3=0{,}12$.
	} 
\end{ex}

% nguồn: Ok_TracNghiemLop11_Lan1
\begin{ex}%[1D9N1-3]%[ma:28492e80]
	Cho $A,B$ là hai biến cố độc lập.  Biết $\mathrm{P}(A)=\dfrac{1}{3}, \mathrm{P}(B)=\dfrac{1}{4}$. Giá trị của $\mathrm{P}(A\cap B)$ bằng
	\choice 
		{\True  $\dfrac{1}{12}$}
		{$\dfrac{1}{7}$} 
		{$\dfrac{7}{12}$} 
		 {$\dfrac{1}{2}$}
	\loigiai{ 
		Vì  $A,B$ là hai biến cố độc lập nên ta có\\
		\hspace*{5cm}$\mathrm{P}(A\cap B)=\mathrm{P}(A)\cdot  \mathrm{P}(B)=\dfrac{1}{3}\cdot\mathrm{P}(B)=\dfrac{1}{12}$.
	 }
\end{ex}

% nguồn: Ok_TracNghiemLop11_Lan1
\begin{ex}%[1D9N1-3]%[ma:d9fcd67a]
	Hai người độc lập nhau ném bóng vào rổ. Mỗi người ném vào rổ của mình một quả bóng. Biết rằng xác suất ném  bóng trúng vào rổ của từng người tương ứng là $\dfrac{1}{3}$ và $\dfrac{3}{7}$. Gọi $A$ là biến cố: \lq\lq Cả hai cùng ném bóng trúng vào rổ \rq\rq. Khi đó, xác suất của biến cố $A$ là bao nhiêu?
	 \choice
		{$\dfrac{1}{3}$} 
		 {\True $\dfrac{1}{7}$}
		{$\dfrac{1}{21}$} 
		 {$\dfrac{3}{7}$}
	 \loigiai{Do hai người ném bóng vào rổ là độc lập nên theo quy tắc nhân xác suất của hai biến cố độc lập suy ra $P(A)=\dfrac{1}{3}\cdot \dfrac{3}{7}=\dfrac{1}{7}.$
	} 
\end{ex}

% nguồn: Ok_TracNghiemLop11_Lan1
\begin{ex}%[1D9N1-3]%[ma:a1eee090]
	Cho $A$, $B$ là  hai biến cố độc lập. Biết $\mathrm{P}(A) = \dfrac{1}{4}$,$\mathrm{P}(AB) = \dfrac{1}{9}$. Tính $\mathrm{P}(B)$.
	 \choice
		 {\True $\dfrac{4}{9}$}
		{$\dfrac{5}{36}$} 
		 {$\dfrac{13}{36}$}
		 {$\dfrac{9}{4}$}
	 \loigiai{
		Ta có  $\mathrm{P}(AB) = \mathrm{P}(A) \cdot \mathrm{P}(B) \Leftrightarrow \mathrm{P}(B) = \dfrac{\mathrm{P}(AB)}{\mathrm{P}(A)} = \dfrac{4}{9}$.
	} 
\end{ex}

% nguồn: Ok_TracNghiemLop11_Lan1
\begin{ex}%[1D9N1-3]%[ma:01b3c6a4]
	Cho $A$ và $B$ là hai biến cố độc  lập với nhau, $\mathrm{P}(A)=0{,}4$ và $\mathrm{P}(AB)=0{,}02$. Khi đó $\mathrm{P}(B)$ bằng
	 \choice
		 {$0{,}55$}
		{$0{,}06$} 
		{\True $0{,}05$} 
		 {$0{,}25$}
	\loigiai{ 
		Ta có $A$ và $B$ là hai biến cố độc lập với  nhau, do đó $\mathrm{P}(AB)=\mathrm{P}(A)\cdot \mathrm{P}(B) \Rightarrow \mathrm{P}(B)=\dfrac{\mathrm{P}(AB)}{\mathrm{P}(A)}=\dfrac{0{,02}}{0{,}4}=0{,}05$.
	 }
\end{ex}

% nguồn: Ok_TracNghiemLop11_Lan1
\begin{ex}%[1D9H1-3]%[ma:e375a446]
	Hai xạ thủ mỗi người bắn một viên đạn vào bia một cách độc lập. Xác suất bắn trúng bia của xạ thủ thứ nhất và xạ thủ thứ hai lần lượt là $0{,}9$ và $0{,}7$. Xác suất để người thứ nhất bắn trúng và người  thứ hai bắn không trúng bia là
	 \choice
		 {\True $0{,}27$}
		 {$0{,}63$}
		{$0{,}07$} 
		 {$0{,}72$}
	 \loigiai{
		Xác suất để người thứ nhất bắn trúng  và người thứ hai bắn không trúng bia là $0{,}9\cdot 0{,}3=0{,}27$.
	 }
\end{ex}

% nguồn: Ok_TracNghiemLop11_Lan1
\begin{ex}%[1D9H1-3]%[ma:2837b482]
	Một hộp chứa $22$ tấm thẻ cùng loại được đánh số lần lượt từ $1$ đến $22$. Chọn ra ngẫu nhiên $1$ thẻ từ hộp. Gọi $A$ là biến cố \lq\lq Số ghi trên thẻ được chọn chia hết cho $2$\rq\rq,  $B$ là biến cố \lq\lq Số ghi trên thẻ được chọn chia hết cho $3$\rq\rq. Xác suất của biến cố $AB$ bằng
	 \choice
		{\True  $\dfrac{3}{22}$}
		 {$\dfrac{7}{22}$}
		{$\dfrac{1}{2}$} 
		{$\dfrac{7}{44}$} 
	 \loigiai{
		Không gian mẫu $\Omega=\left\{1;\; 2;\; 3;\;...;\; 22\right\}$,  suy ra $n\left(\Omega \right)=22$.\\
		Ta  có $A$ là biến cố \lq\lq Số ghi trên thẻ được chọn chia hết cho $2$\rq\rq \, $\Rightarrow A=\left\{2;\; 4;\; 6;\;...;\; 22\right\}$.\\
		$B$ là biến cố  \lq\lq Số ghi trên thẻ được chọn chia hết cho $3$\rq\rq \, $\Rightarrow B=\left\{3;\; 6;\; 9;\; 12;\; 15;\; 18;\; 21\right\}$.\\
		 Ta cũng có $A\cap B=\left\{6;\; 12;\; 18\right\}$, suy ra $\mathrm{P}(AB)=\dfrac{3}{22}$.
	} 
\end{ex}

% nguồn: Ok_TracNghiemLop11_Lan1
\begin{ex}%[1D9H1-3]%[ma:b5392bf2]
	Một học sinh tô ngẫu nhiên $5$ câu trắc nghiệm, mỗi câu có $4$ phương án trả lời, trong đó có một phương án đúng. Xác suất để học sinh đó tô sai cả $5$ câu là bao nhiêu? 
	\choice 
		{$\dfrac{15}{1024}$} 
		{$\dfrac{3}{4}$} 
		{\True $\dfrac{243}{1024}$} 
		 {$\dfrac{1}{1024}$}
	 \loigiai{
		Gọi $A_i$ là biến cố  \lq\lq Học sinh tô sai ở câu thứ $i$\rq\rq, với $i=\overline{1,\ldots, 5}$.\\
		Ta có $A_1$, $A_2$, $\ldots$, $A_5$ là các biến cố độc lập với $\mathrm{P}(A_i)=\dfrac{3}{4}$.\\ 
		Vậy $\mathrm{P}\left(A_1 A_2 A_3 A_4 A_5\right) = \mathrm{P}(A_1) \cdot  \mathrm{P}(A_2) \cdot \mathrm{P}(A_3) \cdot \mathrm{P}(A_4) \cdot \mathrm{P}(A_5) =\left(\dfrac{3}{4}\right)^5 =\dfrac{243}{1024}$.
	 }
\end{ex}

% nguồn: Ok_TracNghiemLop11_Lan1
\begin{ex}%[1D9H1-3]%[ma:ec589468]
	 Xác suất bắn trúng mục tiêu của một vận động viên khi bắn một viên đạn là $0{,}6$. Người đó bắn hai viên đạn một cách độc lập. Xác suất để một viên trúng mục tiêu và một viên trượt mục tiêu là
	 \choice
		{$0,4$} 
		 {$0,6$}
		{\True $0,48$} 
		 {$0,24$}
	\loigiai{ 
		Gọi $A$ là biến cố  viên đạn thứ nhất bắn trúng mục tiêu.\\
		$B$ là biến cố viên đạn thứ  hai bắn trúng mục tiêu.\\
		$C$ là biến cố có đúng  một viên đạn bắn trúng mục tiêu.\\
		Vì $A$, $B$ là hai biến cố  độc lập suy ra $C=A\cdot \overline{B}\cup\overline{A}\cdot B$.\\
	Suy ra  $\mathrm{P}(C)=\mathrm{P}(A)\cdot \mathrm{P}(\overline{B})+\mathrm{P}(\overline{A})\cdot \mathrm{P}(B)$ $=0{,}6\cdot0{,}4+0{,}6\cdot0{,}4=0{,}48$.}
\end{ex}

% nguồn: Ok_TracNghiemLop11_Lan1
\begin{ex}%[1D9H1-3]%[ma:8616922d]
	Cho $A$ và $B$ là hai biến cố độc lập. Biết $\mathrm{P}(B)=0{,}5$, $\mathrm{P}(A\cup B)=0{,}7$. Tính $\mathrm{P}(A)$ 
	\choice 
		{$0{,}2$} 
		{$0{,}5$} 
		{$0{,}7$} 
		{\True  $0{,}4$}
	 \loigiai{
		 Ta có $\mathrm{P}(A\cup B)=\mathrm{P}(A)+\mathrm{P}(B)-\mathrm{P}(AB)$.\\
		 Vì $A$ và $B$ là các biến cố độc lập nên
		\begin{eqnarray*}
			&  & \mathrm{P}(A\cup B)=\mathrm{P}(A)+\mathrm{P}(B)-\mathrm{P}(A)\cdot \mathrm{P}(B)\\
			&\Leftrightarrow & 0{,}7=\mathrm{P}(A)+0{,}5-0{,}5\cdot \mathrm{P}(A)\\ 
			&\Leftrightarrow &  \mathrm{P}(A)=0{,}4.
		\end{eqnarray*}
	} 
\end{ex}

% nguồn: Ok_TracNghiemLop11_Lan1
\begin{ex}%[1D9H1-3]%[ma:3c69c2e2]
	Một người có một chùm chìa khóa gồm $9$ chiếc, bề ngoài chúng giống hệt nhau và chỉ có đúng hai chiếc mở được cửa nhà. Người đó thử ngẫu  nhiên từng chìa. Xác suất để mở được cửa trong lần mở thứ ba bằng
	\choice 
		 {\True $\dfrac{1}{6}$}
		 {$\dfrac{2}{7}$}
		 {$\dfrac{14}{81}$}
		{$\dfrac{7}{81}$} 
	\loigiai{ 
		Xác suất để mở được cửa trong lần mở thứ  ba là
		 \[
		 \mathrm{P}(A)=\dfrac{7}{9}\cdot \dfrac{6}{8}\cdot \dfrac{2}{7}=\dfrac{1}{6}.
		 \]
	} 
\end{ex}

% nguồn: Ok_TracNghiemLop11_Lan1
\begin{ex}%[1D9H1-3]%[ma:d35f05c6]
	Hai chuyến bay $A$ và $B$ hoạt động độc lập với nhau. Xác suất để chuyến bay $A$ và $B$ bay đúng giờ $0{,}6$ và $0{,}7$. Xác suất để trong $2$ chuyến bay có ít nhất  $1$ chuyến bay đúng giờ bằng
	 \choice
		{\True  $0{,}88$}
		{$0{,}42$} 
		{$0{,}9$} 
		 {$0{,}46$}
	\loigiai{ 
		Xác suất $A$ và  $B$ bay không đúng giờ là $0{,}4$ và $0{,}3$.\\
		Xác suất  cả 2 không đúng giờ là $0{,}4\cdot 0{,}3=0{,}12$.\\
		Xác suất có ít nhất 1 chuyến  bay đúng giờ là $1-0{,}12=0{,}88$.
	 }
\end{ex}

% nguồn: Ok_TracNghiemLop11_Lan1
\begin{ex}%[1D9H1-3]%[ma:3c5ad5b1]
	Ba người cùng bắn vào $1$ bia. Xác suất để người thứ nhất, thứ hai,thứ ba bắn trúng đích lần lượt là $0{,}8$; $0{,}6$; $0{,}5$. Xác suất  để có đúng $2$ người bắn trúng đích bằng
	\choice 
		 {$0{,}24$}
		 {$0{,}96$}
		{\True $0{,}46$} 
		 {$0{,}92$}
	 \loigiai{
		Xác suất để người thứ nhất, thứ hai, thứ ba  bắn trúng đích lần lượt là $\mathrm{P}(A_1)=0{,}8$; $\mathrm{P}(A_2)=0{,}6$; $\mathrm{P}(A_2)=0{,}5$.\\
		Xác suất để có đúng hai người bán trúng đích  bằng
	$$\mathrm{P}(A_1)\cdot \mathrm{P}(A_2)\cdot \mathrm{P}(\overline{A_3})+\mathrm{P}(A_1)\cdot \mathrm{P}(\overline{A_2})\cdot \mathrm{P}(A_3)+\mathrm{P}(\overline{A_1})\cdot \mathrm{P}(A_2)\cdot \mathrm{P}(A_3)=0{,}46.$$} 
\end{ex}

% nguồn: Ok_TracNghiemLop11_Lan1
\begin{ex}%[1D9H1-3]%[ma:dab6dbeb]
	Hai bạn Sơn và Tùng độc, mỗi người tung một con xúc xắc một cách độc lập với nhau. Xác  suất để xúc xắc của bạn Sơn xuất hiện số lẻ, xúc xắc của bạn Tùng xuất hiện số lớn hơn $4$ là
	 \choice
		{\True $\dfrac{1}{6}$} 
		 {$\dfrac{1}{5}$}
		{$\dfrac{1}{7}$} 
		 {$\dfrac{2}{11}$}
	 \loigiai{
		Gọi $A$ là biến cố: \lq\lq xúc xắc của bạn  Sơn xuất hiện số lẻ\rq\rq,\\
		$B$ là biến cố: \lq\lq xúc xắc của bạn Tùng xuất hiện số lớn  hơn $4$\rq\rq.\\
		$AB$ là biến cố: \lq\lq xúc xắc của bạn Sơn  xuất hiện số lẻ và xúc xắc của bạn Tùng xuất hiện số lớn hơn 4\rq\rq.\\
		 Ta có $P(A)=\dfrac{n(A)}{n(\Omega)}=\dfrac{3}{6}=\dfrac{1}{2}; P(B)=\dfrac{n(B)}{n(\Omega)}=\dfrac{2}{6}=\dfrac{1}{3}$.\\
		Xác suất để xúc xắc của bạn Sơn xuất hiện số lẻ, xúc xắc của bạn Tùng xuất hiện số lớn hơn  $4$ là
		$$P(A B)=P(A)\cdot  P(B)=\dfrac{1}{6}.$$
	} 
\end{ex}

% nguồn: Ok_TracNghiemLop11_Lan1
\begin{ex}%[1D9H1-3]%[ma:804ddd7f]
	Gieo một đồng tiền cân đối và đồng chất liên tiếp $3$ lần. Tính xác suất của biến cố $A$ \lq\lq Ít nhất một lần xuất hiện mặt sấp\rq\rq. 
	 \choice
		{$\mathrm{P}(A)=\dfrac{1}{2}$} 
		 {$\mathrm{P}(A)=\dfrac{3}{8}$}
		{\True $\mathrm{P}(A)=\dfrac{7}{8}$} 
		{$\mathrm{P}(A)=\dfrac{1}{4}$} 
	 \loigiai{
		Ta có $\overline{A}$ \lq\lq Không  có lần nào xuất hiện mặt sấp\rq\rq \, hay cả $3$ lần đều mặt ngửa.\\
		Theo quy tắc nhân  xác suất ta có $\mathrm{P}\overline{A})=\dfrac{1}{2}\cdot \dfrac{1}{2}\cdot \dfrac{1}{2}=\dfrac{1}{8}$. \\ Vậy $\mathrm{P}(A)=1-\mathrm{P}\overline{A})=1-\dfrac{1}{8}=\dfrac{7}{8}$.
	 }
\end{ex}

% nguồn: Ok_TracNghiemLop11_Lan1
\begin{ex}%[1D9H1-3]%[ma:54b1fba1]
	Cho $A$ và $B$ là hai biến cố độc lập. Biết $\mathrm{P}(A)=0{,}7$ và $\mathrm{P}(B)=0{,}2$. Xác suất  của biến cố $\overline{A}B$ bằng
	\choice 
		{\True $0{,}06$} 
		{$0{,}9$} 
		 {$0{,}14$}
		 {$0{,}24$}
	 \loigiai{
		Vì $\overline{A}$ là biến cố đối của $A$  nên $\mathrm{P}\left(\overline{A}\right)=1-\mathrm{P}(A)=0{,}3$.\\
		Do $\overline{A}$ và  $B$ độc lập nên $\mathrm{P}\left(\overline{A}B\right)=\mathrm{P}\left(\overline{A}\right)\cdot \mathrm{P}(B)=0{,}3\cdot 0{,}2=0{,}06$.
	 }
\end{ex}

% nguồn: Ok_TracNghiemLop11_Lan1
\begin{ex}%[1D9H1-3]%[ma:d75a9872]
	Hai vận động viên $A$ và $B$ cùng ném bóng vào rổ một cách độc lập với nhau. Xác suất ném bóng trúng vào rổ của hai vận động viên $A$ và $B$ lần  lượt là $\dfrac{1}{5}$ và $\dfrac{2}{7}$. Xác suất của biến cố \lq\lq Cả hai cùng ném bóng trúng vào rổ \rq\rq \, bằng
	\choice 
		{\True  $\dfrac{2}{35}$}
		 {$\dfrac{1}{35}$}
		{$\dfrac{6}{35}$} 
		{$\dfrac{2}{7}$} 
	 \loigiai{
		Do hai người ném bóng vào rổ một cách độc lập với nhau nên xác suất của  biến cố \lq\lq Cả hai cùng ném bóng trúng vào rổ\rq\rq \, là
		\[ 
		 \mathrm{P}=\dfrac{1}{5}\cdot \dfrac{2}{7}=\dfrac{2}{35}.
		\] 
	 }
\end{ex}

% nguồn: Ok_TracNghiemLop11_Lan1
\begin{ex}%[1D9H1-3]%[ma:e0b8bf58]
	Hai đối thủ ngang tài nhau, cùng thi đấu với nhau để tranh chức vô địch. Người thắng cuộc là người đầu tiên thắng được $6$ ván đấu. Hết buổi sáng, người I đã thắng $5$ ván, còn người II chỉ mới thắng $3$ ván. Buổi chiều hai người sẽ tiếp tục thi đấu. Xác suất để người I vô địch  bằng
	\choice 
		 {$\dfrac{5}{8}$}
		{$\dfrac{1}{2}$} 
		 {$\dfrac{3}{4}$}
		{\True $\dfrac{7}{8}$} 
	\loigiai{ 
		Xét biến cố người I không vô địch xảy ra khi người  II thắng liên tiếp ba ván buổi chiều.\\
		Xác suất  là $\dfrac{1}{2}\cdot \dfrac{1}{2}\cdot \dfrac{1}{2}=\dfrac{1}{8}$.\\
		Vậy xác suất người I vô địch là  $1-\dfrac{1}{8}=\dfrac{7}{8}$.
	} 
\end{ex}

% nguồn: Ok_TracNghiemLop11_Lan1
\begin{ex}%[1D9H1-3]%[ma:04a6cbe1]
	Một chiếc máy có hai chiếc động cơ I và II chạy độc lập nhau. Xác suất để động cơ I và II chạy tốt lần lượt là $0{,}8$ và $0{,}7$. Xác suất để cả hai động cơ chạy tốt là 
	 \choice
		 {$0{,}24$}
		 {$0{,}94$}
		{$0{,}14$} 
		 {\True $0{,}56$}
	 \loigiai{
		Xác suất để cả hai động cơ đều chạy tốt  là $0{,}8\cdot 0{,}7=0{,}56$.
	} 
\end{ex}

% nguồn: Ok_TracNghiemLop11_Lan1
\begin{ex}%[1D9H1-3]%[ma:7e2b1f5e]
	 Hai học sinh cùng giải $1$ bài tập toán một cách độc lập. Xác suất giải đúng bài tập toán của $2$ học sinh lần lượt là $\dfrac{1}{2}$ và $\dfrac{1}{3}$. Xác suất để cả hai học sinh không giải được bài tập toán bằng
	\choice 
		{$\dfrac{1}{2}$} 
		{$\dfrac{1}{6}$} 
		{\True  $\dfrac{1}{3}$}
		{$\dfrac{5}{6}$} 
	\loigiai{ 
		Xác suất để cả hai  học sinh không giải được bài tập toán bằng $\left(1-\dfrac{1}{2} \right)\left(1-\dfrac{1}{3} \right)=\dfrac{1}{3}$.
	 }
\end{ex}

% nguồn: Ok_TracNghiemLop11_Lan1
\begin{ex}%[1D9H1-3]%[ma:97a8a4d0]
	Ba xạ thủ cùng bắn vào một bia. Xác  suất trúng đích lần lượt là $0{,}5$; $0{,}6$ và $0{,}8$. Xác suất để ít nhất một người bắn trúng bia là
	 \choice
		 {$0{,}24$}
		{$0{,}16$} 
		 {$0{,}82$}
		 {\True $0{,}96$}
	\loigiai{ 
		Xác  suất bắn trượt của ba xạ thủ lần lượt là $0{,}5$; $0{,}4$ và $0{,}2$.\\
		Gọi biến  cố $A\colon$ \lq\lq Ít nhất một người bắn trúng bia\rq\rq.\\
		Suy ra  $\overline{A}\colon$ \lq\lq Cả ba người cùng trượt\rq\rq.\\
		Xác suất cả ba người cùng trượt là $\mathrm{P}(\overline{A})=0{,}5\cdot 0{,}4\cdot 0{,}2=0{,}04$.\\ 
	Vậy xác suất để ít nhất một  người bắn trúng bia là $\mathrm{P}(A)=1-\mathrm{P}(\overline{A})=0{,}96$.}
\end{ex}

% nguồn: Ok_TracNghiemLop11_Lan1
\begin{ex}%[1D9H1-3]%[ma:937e2548]
	Một câu trắc nghiệm có $4$ phương án độc lập, trong đó có $1$ phương án đúng. Một học sinh tô ngẫu nhiên $5$  câu trắc nghiệm. Xác suất để học sinh đó tô sai cả $5$ câu bằng
	\choice 
		{$\dfrac{15}{1024}$} 
		{$\dfrac{3}{4}$} 
		{\True  $\dfrac{243}{1024}$}
		 {$\dfrac{1}{1024}$}
	\loigiai{ 
		Xác suất tô sai $1$ câu  là $\dfrac{3}{4}$.\\
		Vậy xác suất để học sinh đó  tô sai cả $5$ câu $\left(\dfrac{3}{4} \right)^5=\dfrac{243}{1024}$.
	} 
\end{ex}

% nguồn: Ok_TracNghiemLop11_Lan1
\begin{ex}%[1D9H1-3]%[ma:2919428a]
	Hai bệnh nhân M và N bị  nhiễm vi rút cúm H5N1. Biết rằng xác suất bị biến chứng nặng của bệnh nhân M là $0{,}05$ và của bệnh nhân N là $0{,}1$. Khả năng bị biến chứng nặng của hai bệnh nhân là độc lập. Tính xác suất của biến cố \lq\lq Bệnh nhân M bị biến chứng nặng, bệnh nhân N không bị biến chứng nặng\rq\rq.
	 \choice
		{\True  $0{,}045$}
		 {$0{,}005$}
		 {$0{,}095$}
		 {$0{,}855$}
	 \loigiai{
		Gọi $A$ là biến cố \lq\lq Bệnh nhân M bị biến chứng nặng\rq\rq. Ta  có $\mathrm{P}(A)=0{,}05$ và $P(\overline{A})=0{,}95$.\\
		Gọi $B$ là biến cố \lq\lq Bệnh nhân N bị biến  chứng nặng\rq\rq. Ta có $\mathrm{P}(B)=0{,}1$ và $P(\overline{B})=0{,}9$.\\
		Do $A$ và $\overline{B}$ độc  lập nên xác suất bệnh nhân M bị biến chứng nặng, bệnh nhân N không bị biến chứng nặng là $$\mathrm{P}(A\overline{B})=\mathrm{P}(A)\cdot\mathrm{P}(\overline{B})=0{,}05\cdot 0{,}9=0{,}045.$$
	} 
\end{ex}

% nguồn: Ok_TracNghiemLop11_Lan1
\begin{ex}%[1D9H1-3]%[ma:2e3e1c7d]
	Xét phép thử \lq\lq Bạn thứ nhất gieo đồng xu, sau đó bạn thứ hai gieo con xúc  xắc\rq\rq. Gọi $A$ là biến cố \lq\lq Đồng xu xuất hiện mặt ngửa\rq\rq, $B$ là biến cố \lq\lq Con xúc xắc xuất hiện mặt $4$ chấm\rq\rq. Tính xác suất của biến cố $A\overline{B}$.
	\choice 
		{$\dfrac{1}{3}$} 
		{$\dfrac{1}{6}$} 
		{$\dfrac{1}{2}$} 
		{\True  $\dfrac{5}{12}$}
	\loigiai{ 
		 Không gian mẫu của phép thử là
		\[\Omega=\left\{(S, 1),(S, 2),(S,  3),(S, 4),(S, 5),(S, 6),(N, 1),(N, 2),(N, 3),(N, 4),(N, 5),(N, 6)\right\}.
		\] 
		Tập hợp mô tả các  biến cố
		\begin{eqnarray*}
			A & = & \left\{(N, 1),(N, 2),(N, 3),(N, 4),(N, 5),(N,  6)\right\};\\
			 B & = &\left\{(S, 4),(N, 4)\right\};\\
			\overline{B}  & = & \left\{(S, 1),(S, 2),(S, 3),(S, 5),(S, 6),(N, 1),(N, 2),(N, 3),(N, 5),(N, 6)\right\}.
		\end{eqnarray*}
		 Từ đó suy ra $A\overline{B}=\left\{(N, 1),(N, 2),(N, 3),(N, 5),(N, 6)\right\}$.\\
		Vậy  xác suất của biến cố $A\overline{B}$ là $\mathrm{P}\left(A\overline{B}\right)=\dfrac{n\left(A\overline{B}\right)}{n\left(\Omega \right)}=\dfrac{5}{12}$.
	} 
\end{ex}

% nguồn: Ok_TracNghiemLop11_Lan1
\begin{ex}%[1D9H1-3]%[ma:9b7a2331]
	Một chiếc máy có hai chiếc động cơ I và II chạy độc lập nhau. Xác suất để động cơ I và II chạy tốt lần  lượt là $0{,}84$ và $0{,}75$. Xác suất để cả hai động cơ chạy tốt bằng
	\choice 
		{$0{,}37$} 
		{$0{,}04$} 
		 {$0{,}12$}
		{\True  $0{,}63$}
	 \loigiai{
		Xác suất để cả hai động cơ đều chạy tốt là $0{,}84\cdot  0{,}75=0{,}63$.
	 }
\end{ex}

% nguồn: Ok_TracNghiemLop11_Lan1
\begin{ex}%[1D9V1-3]%[ma:c4dc65db]
	Trong đề kiểm tra $15$ phút môn Toán của lớp 11A có $20$ câu trắc nghiệm. Mỗi câu trắc nghiệm có $4$ phương án trả lời, trong đó chỉ có một phương án trả lời đúng. Nam giải chắc chắn đúng $10$ câu, $10$  câu còn lại lựa chọn ngẫu nhiên đáp án. Tính xác suất để Nam đạt được đúng $8$ điểm. Biết rằng mỗi câu trả lời đúng được $0{,}5$ điểm, trả lời sai không bị trừ điểm.
	 \choice
		{$\left(\dfrac{1}{4} \right)^6$} 
		{$\left(\dfrac{1}{4}  \right)^{16}$}
		{\True $\mathrm{C}_{10}^6\cdot \left(\dfrac{1}{4} \right)^6\cdot \left(\dfrac{3}{4}  \right)^4$}
		{$\left(\dfrac{1}{4} \right)^{16}  \left(\dfrac{3}{4} \right)^4$}
	 \loigiai{
		Nam giải chắc chắn đúng $10$ câu nên Nam được chắc chắn $5$ điểm.\\ 
		Để Nam đạt được đúng $8$ điểm thì trong $10$ câu còn lại lựa chọn  ngẫu nhiên đáp án phải đúng $6$ câu, sai $4$ câu.\\
		Xác suất khi đánh ngẫu nhiên đúng  một câu trắc nghiệm là $\dfrac{1}{4}$.\\
		Xác suất khi đánh ngẫu nhiên sai  một câu trắc nghiệm là $\dfrac{3}{4}$.\\
		Chọn $6$ câu trắc nghiệm để  đáp đúng từ $10$ câu trắc nghiệm có $\mathrm{C}_{10}^6$ (cách).\\
		Vậy, xác suất để Nam đạt được đúng $8$ điểm là $\mathrm{C}_{10}^6\cdot \left(\dfrac{1}{4} \right)^6\cdot  \left(\dfrac{3}{4} \right)^4$.
	 }
\end{ex}

% nguồn: Ok_TracNghiemLop11_Lan1
\begin{ex}%[1D9V1-3]%[ma:246ef630]
	Gieo ngẫu nhiên một con xúc xắc cân đối và đồng chất ba lần  liên tiếp. Gọi $A$ là biến cố \lq\lq Tổng số chấm xuất hiện trên ba con xúc xắc bằng $8$\rq\rq, $B$ là biến cố \lq\lq Tích số chấm xuất hiện trên ba con xúc xắc không vượt quá $30$\rq\rq. Xác suất của biến cố $AB$ bằng
	\choice 
		{$\dfrac{5}{36}$} 
		{$\dfrac{1}{8}$} 
		{$\dfrac{1}{9}$} 
		 {\True $\dfrac{7}{72}$}
	\loigiai{ 
		Gieo ngẫu nhiên một con xúc xắc cân đối và đồng  chất ba lần liên tiếp thì không gian mẫu là
		$\Omega=\left\{\left(i;\; j;\; k\right)\; \left|\; 1\le i;\; j;\; k\le 6\right. \right\}$, suy ra  $n\left(\Omega \right)=6^3=216$.\\
		Gọi  $A$ là biến cố \lq\lq Tổng số chấm xuất hiện trên ba con xúc xắc bằng $8$\rq\rq.\\
		Ta có  $8=6+1+1=5+2+1=4+3+1=4+2+2=3+3+2$.\\
		Với bộ $\left(6;\; 1;\; 1\right)$ ta có $3$ hoán vị,  với bộ $\left(5; 2; 1\right)$ ta có $3!$ hoán vị.\\
		Tương tự với các  bộ còn lại.\\
		Suy ra số kết quả thuận lợi cho biến cố $A$  là $n(A)=3+3!+3!+3+3=21$.\\
		Xét biến cố $B$ \lq\lq Tích số chấm xuất hiện trên ba con xúc xắc  không vượt quá $30$.\\
		Ta thấy mọi phần tử thuộc $A$ đều thuộc  $B$, tức là $A\subset B$ nên $A\cap B=A$.\\
		Vậy  xác suất cần tính là $\mathrm{P}(AB)=\mathrm{P}(A)=\dfrac{21}{216}=\dfrac{7}{72}$.
	} 
\end{ex}

% nguồn: Ok_TracNghiemLop11_Lan1
\begin{ex}%[1D9V1-3]%[ma:3c1fd022]
	Một  xạ thủ bắn liên tiếp ba mũi tên vào bia ngắm. Biết rằng xác suất bắn trúng bia ngắm của xạ thủ là $ 0{,}6$. Xác suất để trong ba lần bắn tên có đúng hai lần tên trúng bia ngắm là
	\choice 
		 {$\dfrac{81}{125}$}
		{\True  $\dfrac{54}{125}$}
		{$\dfrac{24}{125}$} 
		{$\dfrac{18}{125}$} 
	 \loigiai{
		$A_1$ là biến cố \lq \lq Mũi tên  thứ nhất của xạ thủ bắn trúng bia\rq \rq.\\
		$A_2$ là biến cố \lq \lq Mũi tên thứ hai của xạ thủ  bắn trúng bia\rq \rq\\
		$A_3$ là biến cố \lq \lq Mũi tên thứ ba của xạ thủ bắn trúng bia\rq \rq.\\ 
		Khi đó, $\mathrm{P}\left(A_1\right)=\mathrm{P}\left(A_2\right)=\mathrm{P}\left(A_3\right)=0{,}6$ và $  \mathrm{P}\left(\overline{A_1}\right)=\mathrm{P}\left(\overline{A_2}\right)=\mathrm{P}\left(\overline{A_3}\right)=0{,}4$.\\
		Gọi $ X$ là biến cố \lq \lq Trong ba lần bắn tên có đúng hai lần  tên trúng bia ngắm\rq \rq.\\
		Ta  có $\mathrm{P}(X)=\mathrm{P}\left(A_1A_2\overline{A_3}\right)+\mathrm{P}\left(A_1\overline{A_2}{A_3}\right)+\mathrm{P}\left(\overline{A_1}{A_2}{A_3}\right).$\\
		Rõ ràng các biến cố $A_1$, $A_2$, $A_3$  đôi một độc lập với nhau nên
		\allowdisplaybreaks 
		\begin{eqnarray*}
			 \mathrm{P}(X)&=&\mathrm{P}\left(A_1\right)\mathrm{P}\left(A_2\right)\mathrm{P}\left(\overline{A_3}\right)+\mathrm{P}\left(A_1\right)\mathrm{P}\left(\overline{A_2}\right)\mathrm{P}\left(A_3\right)+\mathrm{P}\left(A_1\right)\mathrm{P}\left(A_2\right)\mathrm{P}\left(\overline{A_3}\right)\\
			&=&3\cdot 0{,}6\cdot 0{,}6\cdot 0{,}4=\dfrac{54}{125}. 
	\end{eqnarray*} }
\end{ex}


%%%%%%%%%%%%%%%% ĐÚNG/SAI %%%%%%%%%%%%%%%%

% nguồn: Ok_dungsai_lop11_Lan1
\begin{ex}%[1D9H1-3]%[ma:ed00418f]
 Hai bạn An và Hà của lớp 11A tham gia giải bóng bàn đơn nữ do nhà trường tổ chức. Hai bạn đó nằm ở hai bảng đấu loại khác nhau, mỗi bảng đấu loại chỉ chọn một người vào vòng chung kết. Xác suất lọt qua vòng loại để vào vòng  chung kết của An và Hà lần lượt là $0{,}6$ và $0{,}7$. Các khẳng định sau đây, khẳng định nào đúng, khẳng định nào sai?
  \choiceTF
  {\True Biến cố  \lq\lq Bạn An lọt vào vòng chung kết\rq\rq~ và \lq\lq Bạn Hà lọt vào vòng chung kết\rq\rq~ là hai biến cố độc lập}
  {\True Xác suất cả hai bạn  lọt vào vòng chung kết là $0{,}42$}
  {Xác suất có ít nhất một bạn  lọt vào vòng chung kết là $0{,}8$}
  {Xác  suất chỉ có bạn Hà lọt vào vòng chung kết là $0{,}7$}
  \loigiai{
  Xét các  biến cố $A\colon$ \lq\lq Bạn An lọt vào vòng chung kết\rq\rq~ $\Rightarrow \mathrm{P}(A)=0{,}6$.\\
  $B\colon$ \lq\lq Bạn Hà lọt vào vòng  chung kết\rq\rq~ $\Rightarrow \mathrm{P}(B)=0{,}7$.
  \begin{itemchoice}
   \itemch Vì bạn An  và Hà nằm ở hai bảng đấu loại khác nhau nên biến cố $A$ và $B$ độc lập.
    \itemch Ta có $\mathrm{P}(A\cap B)=\mathrm{P}(A)\cdot\mathrm{P}(B)=0{,}6\cdot 0{,}7=0{,}42$.
    \itemch Ta có $\mathrm{P}(A\cup B)=\mathrm{P}(A)+\mathrm{P}(B)-\mathrm{P}(A\cap B)=0{,}6+0{,}7-0{,}42=0{,}88$.
   \itemch Xác suất bạn  An không lọt vào vòng chung kết là $\mathrm{P}(\overline{A})=1-\mathrm{P}(A)=1-0{,}6=0{,}4$.\\
   Xác suất chỉ có bạn Hà lọt vào vòng chung kết là $\mathrm{P}(\overline{A})\cdot\mathrm{P}(B)=0{,}4\cdot  0{,}7=0{,}28$.
  \end{itemchoice}
 } 
\end{ex}

% nguồn: Ok_dungsai_lop11_Lan1
\begin{ex}%[1D9H1-3]%[ma:53f06496]
 Có hai hộp đựng các viên bi. Hộp thứ nhất chứa $5$ viên bi đỏ, $6$ viên bi xanh và $7$ viên bi vàng. Hộp thứ hai chứa $3$ viên  bi trắng, $4$ viên bi đen (các viên bi kích thước như nhau). Các khẳng định sau đúng hay sai?
  \choiceTF
  {Xác suất lấy $1$  viên bi đỏ từ hộp thứ nhất là $\dfrac{1}{18}$}
  {\True Xác suất lấy  $2$ viên bi cùng màu từ hộp thứ hai là $\dfrac{3}{7}$}
  {Xác suất  lấy $4$ viên bi có nhiều nhất $2$ viên bi đỏ từ hộp thứ nhất là $\dfrac{20}{51}$}
  {\True Lấy ngẫu nhiên $4$ viên bi từ hộp thứ nhất và $2$ viên bi từ hộp thứ hai. Xác suất  để trong $6$ viên bi lấy được $1$ viên bi trắng và có ít nhất $3$ viên bi đỏ là $\dfrac{743}{4284}$}
 \loigiai{ 
  \begin{itemchoice}
    \itemch Gọi $A$ là biến cố \lq\lq lấy được một viên bi đỏ từ hộp thứ nhất\rq\rq.\\
    Ta có $n \left(A\right)=5$, $n \left(\Omega\right)= 18 \Rightarrow \mathrm{P}\left(A \right)=\dfrac{5}{18}$.
   \itemch Gọi  $B$ là biến cố \lq\lq lấy được hai viên bi cùng màu từ hộp thứ hai\rq\rq.\\
   Ta có $n  \left(B\right)=\mathrm{C}_3^2+\mathrm{C}_4^2$, $n \left(\Omega\right)= \mathrm{C_7^2} \Rightarrow \mathrm{P}\left(B \right)=\dfrac{3}{7}$.
   \itemch Lấy ngẫu nhiên $4$ viên bi từ hộp  thứ nhất. Số kết quả có thể xảy ra là $n\left(\Omega\right)=\mathrm{C}_{18}^4$.\\
   Gọi $C$ là biến cố \lq\lq $4$ viên bi lấy được trong  hộp thứ nhất có nhiều nhất $2$ viên bi đỏ\rq\rq.\\
   Số kết quả thuận lợi cho biến cố $C$ là  $n\left(C \right)=\mathrm{C}_5^0\cdot \mathrm{C}_{13}^4+\mathrm{C}_5^1\cdot \mathrm{C}_{13}^3+\mathrm{C}_5^2\cdot \mathrm{C}_{13}^2$.\\
   Khi đó, xác suất của  biến cố $C$ là $\mathrm{P}\left(C \right)=\dfrac{n\left(C \right)}{n\left(\Omega\right)}=\dfrac{481}{1224}$.
   \itemch Lấy ngẫu nhiên $2$  viên bi từ hộp thứ hai. Số kết quả có thể xảy ra là $n\left(\Omega\right)=\mathrm{C}_7^2$.\\
   Gọi $D$ là biến cố \lq\lq $2$ viên bi lấy được trong hộp thứ hai  có $1$ viên bi trắng\rq\rq.\\
   Số kết quả thuận lợi cho $B$ là $n\left(D  \right)=\mathrm{C}_3^1\cdot \mathrm{C}_4^1$.\\
   Khi đó, xác suất của biến  cố $D$ là $\mathrm{P}\left(D \right)=\dfrac{n\left(D \right)}{n\left(\Omega\right)}=\dfrac{2}{7}$.\\
   Gọi $E$ là biến cố \lq\lq $4$ viên  bi lấy được trong hộp thứ nhất có ít nhất $3$ viên bi đỏ\rq\rq.\\
   Ta  có $\mathrm{P}\left(E \right)=\mathrm{P}\left(\overline{C} \right)=1-\mathrm{P}\left(C \right)$.\\
   Vậy xác suất để trong $6$ viên bi lấy được có $1$ viên bi trắng và có ít nhất $3$  viên bi đỏ là
   $$\mathrm{P}\left(DE \right)=\mathrm{P}\left(D \right)\cdot \mathrm{P}\left(E \right)=\mathrm{P}\left(D \right)\cdot \left(1-\mathrm{P}\left(C \right) \right)=\dfrac{2}{7}\cdot \left(1-\dfrac{481}{1224} \right)=\dfrac{743}{4284}.$$ 
  \end{itemchoice}
 } 
\end{ex}

% nguồn: Ok_dungsai_lop11_Lan1
\begin{ex}%[1D9H1-3]%[ma:37d55025]
 Gieo hai  con súc sắc I và II cân đối, đồng chất một cách độc lập. Xét các biến cố $A,\,B,\,C$ sau đây\\
 \indent \indent $A\colon$ \lq\lq Có ít nhất một con súc sắc  xuất hiện mặt $6$ chấm\rq\rq;\\
 \indent \indent $B\colon$ \lq\lq Tổng số chấm trên mặt xuất hiện  của hai con là $7$\rq\rq;\\
 \indent \indent $C\colon$ \lq\lq Tổng số chấm trên mặt xuất  hiện của hai con lớn hơn hoặc bằng $8$\rq\rq.\\
 Trong các khẳng định sau đây, khẳng định nào đúng, khẳng  định nào sai?
  \choiceTF
   {Số phần tử của không gian mẫu bằng $12$}
  {$\mathrm{P}(B)=\dfrac{1}{18}$} 
  {\True Hai biến cố $A$ và $B$  không độc lập}
  {$\mathrm{P}(C)=\dfrac{19}{36}$} 
  \loigiai{
  \begin{itemchoice}
   \itemch Gieo  hai con súc sắc I và II cân đối, đồng chất một cách độc lập thì số phần tử của không gian mẫu là $n(\Omega)=6\cdot 6=36$.
   \itemch Do đó khẳng định $B=\left\{(1;6),(2;5),(3;4),(4;3),(5;2),(6;1) \right\}$, nên  $n(B)=6$.\\
   Do  đó $\mathrm{P}(B)=\dfrac{6}{36}=\dfrac{1}{6}$.
   \itemch Ta có $\overline{A}\colon$ \lq\lq Không có con súc sắc nào  xuất hiện mặt $6$ chấm\rq\rq.\\
   Do đó $n(\overline{A})=5\cdot 5=25 \Rightarrow \mathrm{P}(\overline{A})=\dfrac{25}{36} \Rightarrow \mathrm{P}(A)=1-\dfrac{25}{36}=\dfrac{11}{36}$.\\ 
   Suy ra  $\mathrm{P}(A)\cdot\mathrm{P}(B)=\dfrac{11}{36}\cdot\dfrac{1}{6}=\dfrac{11}{216}$.\\
   Xét biến cố $AB\colon$ \lq\lq Tổng số chấm  trên mặt xuất hiện của hai con là $7$, trong đó có ít nhất một con súc sắc xuất hiện mặt $6$ chấm\rq\rq.\\
   Ta có $AB=\left\{(1;6),(6;1) \right\}$. Do đó $\mathrm{P}(AB)=\dfrac{2}{36}$.\\ 
   Vậy $\mathrm{P}(AB)\ne \mathrm{P}(A)\cdot\mathrm{P}(B)$, suy ra $A,\,B$ không độc lập. 
   \itemch Biến cố $C$ xảy ra khi và  chỉ khi một trong các biến cố sau xảy ra
   \begin{itemize}
    \item $C_1\colon$  \lq\lq Con súc sắc I ra mặt $2$ chấm và con súc sắc II ra mặt $6$ chấm\rq\rq.\\
    Suy  ra $\mathrm{P}\left(C_1 \right)=\dfrac{1}{6}\cdot\dfrac{1}{6}=\dfrac{1}{36}$.
    \item $C_2\colon$ \lq\lq Con súc sắc I ra mặt $3$ chấm và con  súc sắc II ra mặt $5$ hoặc $6$ chấm\rq\rq.\\
    Suy ra $\mathrm{P}\left(C_2  \right)=\dfrac{1}{6}\cdot\dfrac{2}{6}=\dfrac{2}{36}$.
    \item $C_3\colon$ \lq\lq Con súc  sắc I ra mặt $4$ chấm và con súc sắc II ra mặt $4,\,5$ hoặc $6$ chấm\rq\rq.\\
    Suy  ra $\mathrm{P}\left(C_3 \right)=\dfrac{1}{6}\cdot\dfrac{3}{6}=\dfrac{3}{36}$.
    \item $C_4\colon$  \lq\lq Con súc sắc I ra mặt $5$ chấm và con súc sắc II ra mặt $3,\,4,\,5$ hoặc $6$ chấm\rq\rq.\\
    Suy ra  $\mathrm{P}\left(C_4 \right)=\dfrac{1}{6}\cdot\dfrac{4}{6}=\dfrac{4}{36}$.
    \item $C_5\colon$ \lq\lq Con súc sắc I ra mặt $6$ chấm và con súc sắc  II ra mặt $2,\,3,\,4,\,5$ hoặc $6$ chấm\rq\rq.\\
    Suy ra $\mathrm{P}\left(C_5  \right)=\dfrac{1}{6}\cdot\dfrac{5}{6}=\dfrac{5}{36}$.
   \end{itemize}
   Theo quy tắc cộng  sác suất ta được
   $$\mathrm{P}(C)=\mathrm{P}\left(C_1  \right)+\mathrm{P}\left(C_2 \right)+\mathrm{P}\left(C_3 \right)+\mathrm{P}\left(C_4 \right)+\mathrm{P}\left(C_5 \right)=\dfrac{1}{36}+\dfrac{2}{36}+\dfrac{3}{36}+\dfrac{4}{36}+\dfrac{5}{36}=\dfrac{15}{36}.$$
  \end{itemchoice}
 } 
\end{ex}

% nguồn: Ok_dungsai_lop11_Lan1
\begin{ex}%[1D9H1-3]%[ma:4d9a6175]
 Cho $A$ và $B$ là hai biến cố độc lập với nhau, biết $\mathrm{P}(A)=0{,}2$; $\mathrm{P}(B)=0{,}3$. Khi đó 
 \choiceTF 
  {\True $\mathrm{P}(AB)=0{,}06$} 
   {$\mathrm{P}(A\overline{B})=0{,}12$}
   {\True $\mathrm{P}(\overline{A}\overline{B})=0{,}56$}
   {\True $\mathrm{P}(\overline{A}B)=0{,}24$}
 \loigiai{ 
  Vì $A$, $B$ là hai biến cố độc lập nên $A$, $B$, $\overline{A}$, $\overline{B}$ là các  biến cố đôi một độc lập.\\
  Ta có 
  \begin{itemchoice}
   \itemch  Đúng. $\mathrm{P}(AB)=\mathrm{P}(A)\cdot\mathrm{P}(B)=0{,}2\cdot 0{,}3=0{,}06$.
   \itemch Sai.  $\mathrm{P}(A\overline{B})=\mathrm{P}(A)\cdot\mathrm{P}(\overline{B})=0{,}2\cdot 0{,}7=0{,}14$.
   \itemch Đúng.  $\mathrm{P}(\overline{A}\overline{B}) =\mathrm{P}(\overline{A})\cdot\mathrm{P}(\overline{B})=0{,}8\cdot 0{,}7=0{,}56$.
   \itemch  Đúng. $\mathrm{P}(\overline{A}B)
   =\mathrm{P}(\overline{A})\cdot\mathrm{P}(B)=0{,}8\cdot  0{,}3=0{,}24$.
  \end{itemchoice}
 } 
\end{ex}

% nguồn: Ok_dungsai_lop11_Lan1
\begin{ex}%[1D9H1-3]%[ma:629a0baa]
 Cho $A$ và $B$ là hai biến cố độc  lập. Biết $\mathrm{P}(A)=0{,}6$ và $\mathrm{P}(B)=0{,}8$. Khi đó
  \choiceTF
  {\True  $\mathrm{P}(AB)=0{,}48$}
  {\True $\mathrm{P}(\overline{A} B)=0{,}32$} 
   {$\mathrm{P}(A \overline{B})=0{,}012$}
  {\True  $\mathrm{P}(\overline{A}\, \overline{B})=0{,}08$}
  \loigiai{
  \begin{itemchoice}
   \itemch Do $A$ và $B$ là hai biến  cố độc lập nên
   $$\mathrm{P}(A B)=\mathrm{P}(A)  \mathrm{P}(B)=0{,}48.$$
   \itemch Vì $\overline{A}$ là biến cố đối của  $A$ nên $\mathrm{P}(\overline{A})=1-\mathrm{P}(A)=0{,}4$. Do $\overline{A}$ và $B$ độc lập nên
   $$\mathrm{P}(\overline{A}  B)=\mathrm{P}(\overline{A}) \mathrm{P}(B)=0{,}32.$$
   \itemch Vì  $\overline{B}$ là biến cố đối của $B$ nên $\mathrm{P}(\overline{B})=1-\mathrm{P}(B)=0{,}2$. Do $\overline{B}$ và $A$ độc lập nên
   $$\mathrm{P}(\overline{B}  A)=\mathrm{P}(\overline{B}) \mathrm{P}(A)=0{,}12.$$
   \itemch Vì  $\overline{B}$ là biến cố đối của $B$ nên $\mathrm{P}(\overline{B})=1-\mathrm{P}(B)=0{,}2$. Do $\overline{A}$ và $\overline{B}$ độc lập nên
   $$\mathrm{P}(\overline{A}\, \overline{B})=\mathrm{P}(\overline{A}) \mathrm{P}(\overline{B})=0{,}08. $$ 
  \end{itemchoice}
 } 
\end{ex}

% nguồn: Ok_dungsai_lop11_Lan1
\begin{ex}%[1D9H1-3]%[ma:5d402918]
 Chọn ngẫu nhiên một lá bài từ bộ bài tú lơ khơ $52$ lá, trả lại lá bài vừa rút vào bộ bài và rút tiếp một lá khác. Xét biến cố $A\colon$ \lq\lq Lần đầu rút được lá $2$\rq\rq! và $B\colon$  \lq\lq Lần hai rút được lá Át\rq\rq. Hãy chỉ ra các khẳng định đúng và các khẳng định sai?
 \choiceTF 
   {\True Hai biến cố $A$ và $B$ độc lập}
  {\True Xác xuất của biến  cố $A$ bằng $\dfrac{1}{13}$}
  {Xác xuất để lần đầu rút lá $2$ và lần hai rút được lá Át bằng $\dfrac{2}{13}$} 
  {\True Xác xuất để  trong hai lá bài rút ra có đủ hai chất cơ là $\dfrac{1}{16}$}
 \loigiai{ 
  \begin{itemchoice}
   \itemch Ta  có
   \begin{itemize}
    \item  Nếu $A$ xảy ra\\
    Vì lá bài rút ra ở lần đầu lại trả lại nên trong bộ bài vẫn đủ $52$ lá và có $4$  lá át.\\
    Do đó xác xuất của biến cố B là  $\mathrm{P}(B)=\dfrac{\mathrm{C}_4^1}{\mathrm{C}_{52}^1}=\dfrac{1}{13}$.
    \item Nếu $A$ không xảy  ra\\
    Vì lá bài rút ra ở lần đầu lại trả  lại nên trong bộ bài vẫn đủ $52$ lá và có $4$ lá át.\\
    Do đó xác xuất của biến cố  B là $\mathrm{P}(B)=\dfrac{\mathrm{C}_4^1}{\mathrm{C}_{52}^1}=\dfrac{1}{13}$.
   \end{itemize}
   Do đó $A$ xảy ra hay không xảy ra không làm ảnh hưởng đến xác suất của biến  cố $B$.\\
   Ngược lại  việc xảy ra hay không xảy ra của biến cố $B$ không làm ảnh hưởng đến $A$ do $B$ rút sau.\\
   Vậy $A$ và $B$  độc lập.
   \itemch  Xác xuất của biến cố $A$ là $\mathrm{P}(A)=\dfrac{\mathrm{C}_4^1}{\mathrm{C}_{52}^1}=\dfrac{1}{13}$.
   \itemch Vì $A$ và $B$ độc lập nên xác xuất để lần đầu rút lá $2$ và lần hai rút được lá  Át là
   $$\mathrm{P}(AB)=\dfrac{\mathrm{C}_4^1}{\mathrm{C}_{52}^1}\cdot\dfrac{\mathrm{C}_4^1}{\mathrm{C}_{52}^1}=\dfrac{1}{13}\cdot\dfrac{1}{13}=\dfrac{1}{169}.$$ 
   \itemch Vì trong $52$ lá bài có $13$ lá cơ nên xác xuất  để lấy được $2$ lá cơ là $\dfrac{13}{52}\cdot\dfrac{13}{52}=\dfrac{1}{16}$.
  \end{itemchoice}
  }
\end{ex}

% nguồn: Ok_dungsai_lop11_Lan1
\begin{ex}%[1D9H1-3]%[ma:317b71c4]
 Cho $A$ và  $B$ là hai biến cố độc lập với nhau, biết $\mathrm{P}(A)=0{,}2; \mathrm{P}(B)=0{,}3$. Khi đó
  \choiceTF
  {\True $\mathrm{P}(A B)=0{,}06$} 
  {$\mathrm{P}(A \overline{B})=0{,}12$} 
   {\True$\mathrm{P}(\overline{A} \overline{B})=0{,}56$}
  {\True$\mathrm{P}(\overline{A}  B)=0{,}24$}
  \loigiai{
  Vì  $A$, $B$ là hai biến cố độc lập nên $A$, $B$, $\overline{A}$, $\overline{B}$ là các biến cố đôi một độc lập.\\
  Ta có\\ 
  $\mathrm{P}(A  B)=P(A) \cdot \mathrm{P}(B)=0{,}2 \cdot 0,3=0{,}06$;
   $\mathrm{P}(A \overline{B})=P(A) \cdot \mathrm{P}(\overline{B})=0{,}2 \cdot 0{,}7=0{,}14$; \\
  $\mathrm{P}(\overline{A} \overline{B})=\mathrm{P}(\overline{A}) \cdot \mathrm{P}(\overline{B})=0{,}8 \cdot 0,7=0{,}56$; 
  $\mathrm{P}(\overline{A} B)=\mathrm{P}(\overline{A})  \cdot \mathrm{P}(B)=0{,}8\cdot0{,}3=0{,}24$.
 } 
\end{ex}

% nguồn: Ok_dungsai_lop11_Lan1
\begin{ex}%[1D9H1-3]%[ma:c85a3587]
 Hai người cùng đi câu cá trên một hồ. Xác suất câu được cá của người thứ nhất và người thứ hai lần lượt là $0{,}8$ và $0{,}9$. Gọi biến cố $A_i \colon$ \lq\lq Người thứ $i$ câu được cá\rq\rq~ với $i\in \left\{1;2\right\}$. Trong các khẳng định nào sau  đây, khẳng định nào đúng, khẳng định nào sai?
 \choiceTF 
  {\True $A_1;A_2$ là hai biến cố độc  lập}
  {Xác suất biến cố \lq\lq Cả hai người câu được cá \rq\rq~ là  $0{,}96$}
  {\True Xác suất biến cố \lq\lq Người thứ nhất câu được cá, người  thứ hai không câu được cá\rq\rq~ là $0{,}08$}
  {\True Xác suất biến cố \lq\lq Có ít nhất một người câu được cá\rq\rq~ là $0{,}98$} 
 \loigiai{ 
  Ta có biến cố $A_i\colon$ \lq\lq Người thứ $i$ câu  được cá\rq\rq~với $i=1,2$.\\
  Biến cố $\overline{A}_i\colon$ \lq\lq Người thứ $i$ không  câu được cá\rq\rq~với $i=1,2$.\\
   Theo giả thiết ta có $\mathrm{P}\left(A_1 \right)=0{,}8;\mathrm{P}\left(A_2 \right)=0{,}9;\mathrm{P}\left(\overline{A}_1 \right)=0{,}2;\mathrm{P}\left(\overline{A}_2 \right)=0{,}1$.
  \begin{itemchoice}
   \itemch Vì hai người câu khác nhau nên $A_1;A_2$ là hai biến cố độc lập. 
   \itemch Gọi biến cố $A\colon$ \lq\lq Cả hai người câu được cá\rq\rq.\\ 
   Ta có $A=A_1A_2$ và $A_1;A_2$  là hai biến cố độc lập $\Rightarrow \mathrm{P}\left(A \right)=\mathrm{P}\left(A_1 \right)\cdot\mathrm{P}\left(A_2 \right)=0{,}8\cdot 0{,}9=0{,}72$.
   \itemch Gọi biến cố $B\colon$ \lq\lq Người thứ  nhất câu được cá, người thứ hai không câu được cá\rq\rq.\\
   Ta có $B=A_1\overline{A}_2$ và $A_1;\overline{A}_2$ là hai biến cố độc lập $\Rightarrow \mathrm{P}\left(B  \right)=\mathrm{P}\left(A_1 \right)\cdot\mathrm{P}\left(\overline{A}_2 \right)=0{,}8\cdot 0{,}1=0,08$.
   \itemch Gọi biến  cố $C\colon$ \lq\lq Có ít nhất một người câu được cá\rq\rq.\\
    Ta có biến cố $\overline{A}\colon$ \lq\lq Cả hai người không câu được cá\rq\rq.\\
   Hai biến cố trên  là hai biến cố đối nhau nên
    $$\mathrm{P}\left(C \right)=1-\mathrm{P}\left(\overline{A} \right)=1-\mathrm{P}\left(\overline{A}_1 \right)\cdot\mathrm{P}\left(\overline{A}_2 \right)=1-0{,}2.0{,}1=0{,}98.$$
  \end{itemchoice}
 } 
\end{ex}

% nguồn: Ok_dungsai_lop11_Lan1
\begin{ex}%[1D9H1-3]%[ma:e793703d]
 Hai bệnh nhân $X$ và $Y$ bị nhiễm Vì rút SARS-CoV-2. Biết rằng xác suất bị biến chứng nặng của bệnh nhân $X$ là $0{,}1$ và của bệnh nhân $Y$ là $0{,}2$. Khả năng bị biến chứng nặng của hai bệnh nhân là  độc lập. Khi đó
  \choiceTF
  {\True Xác xuất của biến cố  \lq\lq Cả hai bệnh nhân đều bị biến chứng nặng\rq\rq\, là $0{,}02$}
  {Xác xuất của biến cố \lq\lq Cả hai bệnh nhân đều không bị biến  chứng nặng\rq\rq\, là $0{,}7$}
  {\True Xác xuất của biến cố \lq\lq Bệnh nhân $X$ bị biến chứng nặng,  bệnh nhân $Y$ không bị biến chứng nặng\rq\rq\, là $0{,}08$}
  {Xác xuất của biến cố \lq\lq Bệnh nhân $X$ không bị biến chứng nặng, bệnh nhân $Y$  bị biến chứng nặng\rq\rq\, là $0{,}08$}
  \loigiai{
  Gọi $A$ là biến cố \lq\lq Bệnh nhân $X$ bị biến chứng nặng\rq\rq.  Ta có $\mathrm{P}(A)=0{,}1$ và $\mathrm{P}(\overline{A})=0{,}9$.\\
  Gọi $B$ là biến cố \lq\lq Bệnh nhân $Y$ bị biến chứng nặng\rq\rq. Ta có $\mathrm{P}(B)=0{,}2$  và $\mathrm{P}(\overline{B})=0{,}8$.
  \begin{itemchoice}
   \itemch \lq\lq Cả hai bệnh nhân đều bị biến  chứng nặng\rq\rq.\\
   Ta thấy $A$ và $B$ là hai biến cố độc lập nên xác suất cả hai bệnh  nhân đều bị biến chứng nặng là
   $$ \mathrm{P}(A B)=\mathrm{P}(A) \mathrm{P}(B)=0{,}02.$$ 
   \itemch \lq\lq  Cả hai bệnh nhân đều không bị biến chứng nặng\rq\rq.\\
   Do  $\overline{A}$ và $\overline{B}$ độc lập nên xác suất cả hai bệnh nhân không bị biến chứng nặng là
   $$ \mathrm{P}(\overline{A} \overline{B})=\mathrm{P}(\overline{A}) \mathrm{P}(\overline{B})=0{,}72. 
   $$ 
   \itemch \lq\lq Bệnh nhân $X$ bị biến chứng nặng, bệnh nhân  $Y$ không bị biến chứng nặng\rq\rq.\\
   Do $A$ và $\overline{B}$ độc lập nên  xác suất bệnh nhân $X$ bị biến chứng nặng, bệnh nhân $Y$ không bị biến chứng nặng là
   $$\mathrm{P}(A  \overline{B})=\mathrm{P}(A) \mathrm{P}(\overline{B})=0{,}08.$$
   \itemch  \lq\lq Bệnh nhân $X$ không bị biến chứng nặng, bệnh nhân $Y$ bị biến chứng nặng\rq\rq.\\
   Do  $\overline{A}$ và $B$ độc lập nên xác suất bệnh nhân $X$ không bị biến chứng nặng, bệnh nhân $Y$ bị biến chứng nặng là
   $$\mathrm{P}(B \overline{A})=\mathrm{P}(B)  \mathrm{P}(\overline{A})=0{,}18.$$
  \end{itemchoice}
  }
\end{ex}

% nguồn: Ok_dungsai_lop11_Lan1
\begin{ex}%[1D9H1-3]%[ma:a121d624]
 Một cuộc thi bắn súng, có $3$ người tham gia thi. Trong đó xác suất bắn trúng của người thứ nhất  là $0{,}9$; người thứ $2$ là $0{,}7$ và người thứ $3$ là $0{,}8$. Các khẳng định sau đây, khẳng định nào là đúng, khẳng định nào là sai?
  \choiceTF
  {\True Xác suất để cả ba người đều bắn trúng  là $0{,}504$}
   {\True Xác suất để có đúng hai người bắn trúng là $0{,}398$}
  {Xác suất để không có người nào bắn trúng là $0{,}4$} 
  {\True Xác suất để có ít nhất một người bắn trúng là $0{,}994$} 
 \loigiai{ 
   Gọi $A$ là biến cố \lq\lq Người thứ nhất bắn trúng\rq\rq, $\mathrm{P}(A)=0{,}9$.\\
  $B$  là biến cố \lq\lq Người thứ hai bắn trúng\rq\rq, $\mathrm{P}(B)=0{,}7$.\\
  $C$ là biến cố  \lq\lq Người thứ ba bắn trúng\rq\rq, $\mathrm{P}(C)=0{,}8$.\\
  Suy ra $A,\,B,\,C$  là ba biến cố độc lập. Khi đó\\
  $\overline{A}$ là  biến cố \lq\lq Người thứ nhất không bắn trúng\rq\rq; $\mathrm{P}(\overline{A})=1-0{,}9=0{,}1$.\\
  $\overline{B}$ là biến cố \lq\lq Người thứ hai bắn không  trúng\rq\rq; $\mathrm{P}(\overline{B})=1-0{,}7=0{,}3$.\\
  $\overline{C}$ là biến cố \lq\lq Người thứ ba bắn không  trúng\rq\rq; $\mathrm{P}(\overline{C})=1-0{,}8=0{,}2$.
  \begin{itemchoice}
   \itemch $A\cap B\cap C$ là  biến cố \lq\lq Cả ba người bắn trúng\rq\rq.\\
   Xác suất để cả ba người bắn trúng là $\mathrm{P}(A\cap B\cap C)=\mathrm{P}(A).\mathrm{P}(B)\cdot\mathrm{P}(C)=0{,}9\cdot 0{,}7\cdot 0{,}8=0{,}504$. 
   \itemch Gọi $D$ là biến cố \lq\lq Có đúng hai người bắn trúng\rq\rq. Ta có $D=(A\cap B\cap \overline{C})\cup (A\cap \overline{B}\cap  C)\cup (\overline{A}\cap B\cap C)$.\\
   Xác suất để  có đúng hai người bắn trúng là $\mathrm{P}(D)=0{,}9\cdot 0{,}7\cdot 0{,}2+0{,}9\cdot 0{,}3\cdot 0{,}8+0{,}1\cdot 0{,}7\cdot 0{,}8=0{,}398$.
   \itemch Gọi $E=\overline{A}\cap \overline{B}\cap \overline{C}$ là biến  cố \lq\lq Không người nào bắn trúng\rq\rq.\\
    Xác suất để không người nào bắn trúng là $\mathrm{P}(E)=\mathrm{P}(\overline{A}\cap \overline{B}\cap \overline{C})=\mathrm{P}(\overline{A})\cdot\mathrm{P}(\overline{B})\cdot\mathrm{P}(\overline{C})=0{,}1\cdot 0{,}3\cdot 0{,}2=0{,}006$.
   \itemch Gọi $\overline{E}$ là biến cố \lq\lq Ít nhất  một người bắn trúng\rq\rq.\\
   Xác suất để  có ít nhất một người bắn trúng là $\mathrm{P}(\overline{E})=1-\mathrm{P}(E)=1-0{,}006=0{,}994$.
  \end{itemchoice}
 } 
\end{ex}

% nguồn: Ok_dungsai_lop11_Lan1
\begin{ex}%[1D9H1-3]%[ma:7e89a71a]
 Cho $A,B$ là hai biến  cố độc lập và $\mathrm{P}(A)=\dfrac{1}{4}, \mathrm{P}(B)=\dfrac{1}{3}$. Khi đó
  \choiceTF
  {$\mathrm{P}(AB)=\dfrac{1}{2}$} 
  {$\mathrm{P}(A\overline{B})=\dfrac{1}{16}$} 
  {\True  $\mathrm{P}(\overline{A}\overline{B})=\dfrac{1}{2}$}
  {\True  $\mathrm{P}(\overline{A}B)=\dfrac{1}{4}$}
  \loigiai{
   Ta có $\mathrm{P}(\overline{A})=\dfrac{3}{4}, \mathrm{P}(\overline{B})=\dfrac{2}{3}$.\\
  Vì $A$, $B$ là hai biến cố độc lập nên mỗi cặp biến cố lấy ra từ $A, \overline{B}$; $\overline{A}, B$ và $\overline{A}, \overline{B}$ đều là cặp biến cố độc lập. 
  \begin{itemchoice}
    \itemch \textbf{Sai.}\\
   Ta có $\mathrm{P}(AB)= \mathrm{P}(A) \cdot \mathrm{P}(B) = \dfrac{1}{4}  \cdot \dfrac{1}{3} = \dfrac{1}{12}$.
    \itemch \textbf{Sai.}\\
   Ta có $\mathrm{P}(A\overline{B})= \mathrm{P}(A) \cdot \mathrm{P}(\overline{B}) = \dfrac{1}{4}  \cdot \dfrac{2}{3} = \dfrac{1}{6}$.
   \itemch  \textbf{Đúng.}\\
   Ta có $\mathrm{P}(\overline{A}\overline{B})= \mathrm{P}(\overline{A}) \cdot \mathrm{P}(\overline{B}) = \dfrac{3}{4} \cdot \dfrac{2}{3} =  \dfrac{1}{2}$.
   \itemch  \textbf{Đúng.}\\
   Ta có $\mathrm{P}(\overline{A}B)= \mathrm{P}(\overline{A}) \cdot \mathrm{P}(B) = \dfrac{3}{4}  \cdot \dfrac{1}{3} = \dfrac{1}{4}$.
  \end{itemchoice}
 } 
\end{ex}

% nguồn: Ok_dungsai_lop11_Lan1
\begin{ex}%[1D9H1-3]%[ma:d537a0bb]
 Ba xạ thủ cùng bắn mỗi người một viên đạn vào bia một cách độc lập với nhau. Xác suất bắn trúng bia của xạ thủ thứ nhất, thứ hai, thứ ba lần lượt là $\dfrac{1}{2};\,\dfrac{2}{3};\,\dfrac{3}{4}$. Các khẳng định sau đây, khẳng  định nào đúng, khẳng định nào sai?
 \choiceTF 
  {\True Xác suất cả ba xạ thủ đều bắn trúng bia là $\dfrac{1}{4}$} 
  {Xác suất để duy nhất xạ thủ thứ  nhất bắn trúng bia là $\dfrac{1}{2}$}
  {\True Xác  suất có duy nhất một xạ thủ bắn trúng bia là $\dfrac{1}{4}$}
  {Xác suất  có ít nhất hai xạ thủ bắn trúng bia là $\dfrac{3}{4}$}
  \loigiai{
  Gọi $A,B,C$ lần lượt là biến cố xạ thủ thứ nhất, xạ thủ thứ hai và xạ  thủ thứ ba bắn trúng bia.\\
   Ta có $\mathrm{P}(A)=\dfrac{1}{2};\mathrm{P}(B)=\dfrac{2}{3};\mathrm{P}(C)=\dfrac{3}{4} \Rightarrow \mathrm{P}(\overline{A})=\dfrac{1}{2};\mathrm{P}(\overline{B})=\dfrac{1}{3};\mathrm{P}(\overline{C})=\dfrac{1}{4}$.
  \begin{itemchoice}
   \itemch Gọi $X$ là biến cố  cả ba xạ thủ đều bắn trúng bia ta có $X=ABC$.\\
   Do $A,B,C$ là ba biến  cố độc lập nên $\mathrm{P}(X)=\mathrm{P}(A)\cdot\mathrm{P}(B)\cdot\mathrm{P}(C)=\dfrac{1}{2}\cdot\dfrac{2}{3}\cdot\dfrac{3}{4}=\dfrac{1}{4}$.
   \itemch Gọi  $X_1$ là biến cố duy nhất xạ thủ bắn thứ nhất bắn trúng bia ta có $X_1=A\overline{B}\,\overline{C}$.\\
   Do đó $\mathrm{P}\left(X_1  \right)=\mathrm{P}(A)\cdot\mathrm{P}(\overline{B})\cdot\mathrm{P}(\overline{C})=\dfrac{1}{2}\cdot\dfrac{1}{3}\cdot\dfrac{1}{4}=\dfrac{1}{24}$.\\
   \itemch Gọi $X_2$ là biến cố có duy nhất một xạ thủ bắn trúng bia ta có  $X_2=A\overline{B}\,\overline{C}\cup \overline{A}B\overline{C}\cup \overline{A}\,\overline{B}C$.\\
   Do  đó
   \allowdisplaybreaks 
   \begin{eqnarray*}
    \mathrm{P}\left(X_2  \right) &=& \mathrm{P}(A)\cdot\mathrm{P}(\overline{B})\cdot\mathrm{P}(\overline{C})+\mathrm{P}(\overline{A})\cdot\mathrm{P}(B)\cdot\mathrm{P}(\overline{C})+\mathrm{P}(\overline{A})\cdot\mathrm{P}(\overline{B})\cdot\mathrm{P}(C)\\
    &=&  \dfrac{1}{2}\cdot\dfrac{1}{3}\cdot\dfrac{1}{4}+\dfrac{1}{2}\cdot\dfrac{2}{3}\cdot\dfrac{1}{4}+\dfrac{1}{2}\cdot\dfrac{1}{3}\cdot\dfrac{3}{4}=\dfrac{1}{4}.
   \end{eqnarray*}
   \itemch Gọi $X_3$  là biến cố có ít nhất hai xạ thủ không bắn trúng bia.\\
   Ta có $X_3=AB\overline{C}\cup  \overline{A}BC\cup A\overline{B}C\cup ABC$.\\
    Do đó
    \allowdisplaybreaks
   \begin{eqnarray*}
    \mathrm{P}\left(X_3 \right) &=&  \mathrm{P}(A)\cdot\mathrm{P}(B)\cdot\mathrm{P}(\overline{C})+\mathrm{P}(\overline{A})\cdot\mathrm{P}(B)\cdot\mathrm{P}(C)+\mathrm{P}(A)\cdot\mathrm{P}(\overline{B})\cdot\mathrm{P}(C)+\mathrm{P}(A)\cdot\mathrm{P}(B)\cdot\mathrm{P}(C)\\
    &=& \dfrac{1}{2}\cdot\dfrac{2}{3}\cdot\dfrac{1}{4}+\dfrac{1}{2}\cdot\dfrac{2}{3}\cdot\dfrac{3}{4}+\dfrac{1}{2}\cdot\dfrac{1}{3}\cdot\dfrac{3}{4}+\dfrac{1}{2}\cdot\dfrac{2}{3}\cdot\dfrac{3}{4}=\dfrac{17}{24}. 
   \end{eqnarray*}
  \end{itemchoice}
 } 
\end{ex}

% nguồn: Ok_dungsai_lop11_Lan1
\begin{ex}%[1D9H1-3]%[ma:abc1cdf9]
 Ba người cùng bắn vào $1$ bia. Xác suất bắn trúng đích của người thứ nhất, thứ hai, thứ ba lần lượt là $0{,}7$; $0{,}6$; $0{,}8$. Gọi $A$ là biến cố \lq\lq Người thứ nhất bắn trúng đích\rq\rq. Gọi  $B$ là biến cố \lq\lq Người thứ hai bắn trúng đích\rq\rq. Gọi $C$ là biến cố \lq\lq Người thứ ba bắn trúng đích\rq\rq.
 \choiceTF 
  { $ \mathrm{P}(A)=0{,}7$;  $\mathrm{P}(\overline{A})=0{,}7$}
  {\True  $ \mathrm{P}(B)=0{,}6$; $ \mathrm{P}(\overline{B})=0{,}4$}
  {\True  $ \mathrm{P}(C)=0{,}8$; $ \mathrm{P}(\overline{C})=0{,}2$}
  {\True Xác suất để  có đúng $2$ người bắn trúng đích $0{,}452$}
 \loigiai{ 
  \begin{itemchoice}
   \itemch Gọi $A$ là biến cố \lq\lq Người thứ  nhất bắn trúng đích\rq\rq $\Rightarrow \mathrm{P}(A)=0{,}7$; $ \mathrm{P}(\overline{A})=0{,}3$.
   \itemch Gọi $B$ là  biến cố \lq\lq Người thứ hai bắn trúng đích\rq\rq $\Rightarrow \mathrm{P}(B)=0{,}6$; $ \mathrm{P}(\overline{B})=0{,}4$.
   \itemch Gọi $C$ là biến cố \lq\lq Người thứ ba bắn trúng đích\rq\rq  $\Rightarrow \mathrm{P}(C)=0{,}8$; $ \mathrm{P}(\overline{C})=0{,}2$.
   \itemch Gọi $X$ là biến cố \lq\lq Có đúng $2$ người bắn  trúng đích\rq\rq. Vì $A$, $B$, $C$ là ba biến cố độc lập nên ta có
   \begin{eqnarray*}
     \mathrm{P}(X)
     &-& \mathrm{P}(A B\overline{C})+\mathrm{P}(A\overline{B}C)+(\overline{A}B C)\\
    &= &0{,}7\cdot 0{,}6\cdot 0{,}2+0{,}7\cdot  0{,}4\cdot 0{,}8+0{,}3\cdot 0{,}6\cdot 0{,}8\\
     &= & 0{,}452.
   \end{eqnarray*}
  \end{itemchoice}
  }
\end{ex}

% nguồn: Ok_dungsai_lop11_Lan1
\begin{ex}%[1D9H1-3]%[ma:0bc33d6a]
 Gieo hai đồng xu $A$ và $B$ một cách độc lập. Đồng xu $A$ được chế tạo cân đối. Đồng xu $B$ được chế tạo  \textbf{không} cân đối nên xác suất xuất hiện mặt sấp gấp $3$ lần xác suất xuất hiện mặt ngửa.
 \choiceTF 
   {\True Xác suất để đồng xu $A$ xuất hiện mặt ngửa bằng $\dfrac{1}{2}$}
  {\True Xác suất để đồng xu $B$ xuất hiện mặt ngửa bằng:  $\dfrac{1}{4}$}
  {Xác suất để khi gieo hai đồng xu một lần thì cả hai  đều ngửa bằng $\dfrac{1}{12}$}
  {Xác suất để khi gieo hai đồng xu  hai lần thì cả hai đồng xu đều ngửa bằng $\dfrac{1}{32}$}
 \loigiai{ 
  Gọi  $X$ là biến cố \lq\lq Đồng xu $A$ xuất hiện mặt ngửa\rq\rq.\\
  Gọi $Y$ là biến cố \lq\lq Đồng xu $B$ xuất  hiện mặt ngửa\rq\rq.
  \begin{itemchoice}
   \itemch  \textbf{Đúng.}\\
   Vì đồng xu $A$ chế tạo cân đối nên $\mathrm{P}(X)=\dfrac{1}{2}$. 
    \itemch \textbf{Đúng.}\\
   Vì xác suất xuất hiện mặt sấp của đồng xu $B$ gấp $3$ lần xác  suất xuất hiện mặt ngửa của nó nên $\mathrm{P}(Y)=\dfrac{1}{4}$.
   \itemch \textbf{Sai.}\\ 
   Xác suất khi gieo hai đồng xu một lần thì chúng đều ngửa  là
   \[ 
   \mathrm{P}(XY)=\mathrm{P}(X)\cdot \mathrm{P}(Y)=\dfrac{1}{2} \cdot \dfrac{1}{4}=\dfrac{1}{8}. 
   \] 
   \itemch \textbf{Sai.}\\ 
   Xác suất để trong một lần  gieo cả hai đồng xu đều ngửa là $\dfrac{1}{8}$.\\
   Suy ra xác suất khi gieo hai lần thì cả hai lần đó hai  đồng xu đều ngửa là $
   \dfrac{1}{8} \cdot \dfrac{1}{8}=\dfrac{1}{64}. $ 
  \end{itemchoice}
  }
\end{ex}

% nguồn: Ok_dungsai_lop11_Lan1
\begin{ex}%[1D9H1-3]%[ma:7f0734bd]
 Cho $A$, $B$ là hai  biến cố độc lập và $\mathrm{P}(A)=\dfrac{1}{4}$, $\mathrm{P}(B)=\dfrac{1}{3}$. Khi đó
 \choiceTF 
  {$\mathrm{P}(AB)=\dfrac{1}{2}$} 
   {$\mathrm{P}(A\overline{B})=\dfrac{1}{16}$}
  {\True  $\mathrm{P}(\overline{A}\overline{B})=\dfrac{1}{2}$}
  {\True $\mathrm{P}(\overline{A}B)=\dfrac{1}{4}$} 
  \loigiai{
  Ta có xác suất các biến cố đối:  $\mathrm{P}(\overline{A})=\dfrac{3}{4}$, $\mathrm{P}(\overline{B})=\dfrac{2}{3}$.\\
  Vì $A,B$ là hai biến cố độc lập nên mỗi cặp biến cố lấy ra từ  $A$, $B$, $\overline{A}$, $\overline{B}$ đều là cặp biến cố độc lập. Do vậy
  \begin{itemchoice}
    \itemch Sai.
   $\mathrm{P}(AB)=\mathrm{P}(A)\cdot\mathrm{P}(B) =\dfrac{1}{4}\cdot\dfrac{1}{3}=\dfrac{1}{12}$. 
   \itemch  Sai. $\mathrm{P}(\overline{A}B)=\mathrm{P}(\overline{A})\cdot\mathrm{P}(B)
   =\dfrac{3}{4}\cdot\dfrac{1}{3}=\dfrac{1}{4}$. 
   \itemch Đúng. 
    $\mathrm{P}(A\overline{B})=\mathrm{P}(A)\cdot\mathrm{P}(\overline{B})
    =\dfrac{1}{4}\cdot\dfrac{2}{3}=\dfrac{1}{6}$.
    \itemch Đúng. $\mathrm{P}(\overline{A}\overline{B})
    =\mathrm{P}(\overline{A})\cdot\mathrm{P}(\overline{B})=\dfrac{3}{4}\cdot\dfrac{2}{3}=\dfrac{1}{2}$.
 \end{itemchoice}}
\end{ex}

% nguồn: Ok_dungsai_lop11_Lan1
\begin{ex}%[1D9H1-3]%[ma:bb7896be]
 Một nhà xuất bản phát hành hai cuốn sách $A$ và $B$. Thống kê cho thấy có $60\%$ người mua sách $A$; $70\%$ người mua  sách $B$; $50\%$ người mua cả sách $A$ và sách $B$. Chọn ngẫu nhiên một người mua sách. Gọi $A$ là biến cố \lq\lq người đó mua quyển sách $A$\rq\rq, $B$ là biến cố \lq\lq người đó mua quyển sách $B$\rq\rq. Các khẳng định sau đây, khẳng định nào đúng, khẳng định nào sai?
  \choiceTF
  {\True $\mathrm{P}(A)=0{,}6$  và $\mathrm{P}(B)=0{,}7$}
  {\True Biến cố \lq\lq người mua đó mua cả sách  $A$ và sách $B$\rq\rq~ là biến cố $A\cap B$}
  {Xác suất để người mua đó mua ít nhất một trong hai sách  $A$ hoặc $B$ là $0{,}9$}
  {\True Xác suất để người mua đó không mua cả sách $A$ và  sách $B$ là $0{,}2$}
 \loigiai{ 
  \begin{itemchoice}
   \itemch Theo bài ra  ta có $\mathrm{P}(A)=0{,}6;\mathrm{P}(B)=0{,}7$.
   \itemch Biến cố \lq\lq người  mua đó mua cả sách $A$ và sách $B$\rq\rq~ là biến cố $A\cap B$.
   \itemch Biến cố \lq\lq người mua đó mua ít  nhất một trong hai sách $A$ hoặc $B$\rq\rq~ là biến cố $A\cup B$.\\
   Theo bài  ra ta có $\mathrm{P}(A\cap B)=0{,}5$.\\
   Do đó xác  suất để người mua đó mua ít nhất một trong hai sách $A$ hoặc $B$ là
   $$\mathrm{P}(A\cup  B)=\mathrm{P}(A)+\mathrm{P}(B)-\mathrm{P}(A\cap B)=0{,}6+0{,}7-0{,}5=0{,}8.$$
   Vậy xác suất  để người mua đó mua ít nhất một trong hai sách $A$ hoặc $B$ là $0{,}8$.
   \itemch Gọi $C$ là biến cố \lq\lq Người mua  đó không mua cả sách $A$ và sách $B$\rq\rq.\\
   Suy ra biến cố $\overline{C}$ là \lq\lq Người mua đó mua  ít nhất một trong hai sách $A$ hoặc sách $B$\rq\rq.\\
   Do đó $\overline{C}=A\cup B \Rightarrow \mathrm{P}(\overline{C})=0{,}8 \Rightarrow  \mathrm{P}(C)=0{,}2$.\\
   Vậy xác suất để người mua đó không mua cả sách  $A$ và sách $B$ là $0{,}2$.
  \end{itemchoice}
 } 
\end{ex}

% nguồn: Ok_dungsai_lop11_Lan1
\begin{ex}%[1D9H1-3]%[ma:8e383823]
 Một bộ bài tú lơ khơ có $52$ lá, rút ngẫu nhiên lần lượt $3$ lá, mỗi lần rút $1$ lá,  sau mỗi lần rút ta đều để lại lá bài đó vào bộ. Khi đó:
  \choiceTF
  {\True Xác suất rút là bài thứ nhất là con Át  là $\dfrac{4}{52}$}
  {Xác suất rút là bài thứ hai  là con Át là $\dfrac{3}{52}$}
  {Xác suất rút là bài thứ ba  là con $J$ là $\dfrac{1}{52}$}
  {\True Xác suất để hai lần  đầu rút được lá bài Át và lần thứ ba rút được lá bài J là $\dfrac{1}{2197}$}
  \loigiai{
  Gọi $A$ là biến cố \lq\lq Lần thứ nhất rút được  con Át\rq\rq.\\
   Gọi $B$ là biến cố \lq\lq Lần thứ hai rút được con Át\rq\rq.\\
  Gọi $C$ là biến  cố \lq\lq Lần thứ ba rút được con J\rq\rq.\\
  $\Rightarrow ABC$ là biến cố \lq\lq Hai lần đầu  rút được con Át và lần thứ ba rút được con J\rq\rq.\\
  Các  biến cố $A,B$ và $C$ đôi một độc lập với nhau.
  \begin{itemchoice}
   \itemch \textbf{Đúng.}\\ 
   Xác suất rút là bài thứ nhất là con  Át là $\mathrm{P}(A)=\dfrac{4}{52}$.
   \itemch  \textbf{Sai.}\\
   Xác suất rút là bài thứ hai là con Át là  $\mathrm{P}(B)=\dfrac{4}{52}$.
    \itemch \textbf{Sai.}\\
   Xác  suất rút là bài thứ ba là con $J$ là $\mathrm{P}(C)=\dfrac{4}{52}$.
    \itemch \textbf{Đúng.}\\
   Vậy xác suất cần tính là: $\mathrm{P}(ABC)=\mathrm{P}(A)\cdot \mathrm{P}(B)\cdot \mathrm{P}(C)=\dfrac{4}{52} \cdot \dfrac{4}{52} \cdot  \dfrac{4}{52}=\dfrac{1}{2197}$.
  \end{itemchoice}
  }
\end{ex}

% nguồn: Ok_dungsai_lop11_Lan1
\begin{ex}%[1D9H1-3]%[ma:01cb802b]
 Một xạ thủ bắn lần lượt hai viên đạn vào bia. Xác suất bắn không trúng đích của viên thứ nhất và viên thứ hai  lần lượt là $0{,}2$ và $0{,}3$. Biết rằng kết quả các lần bắn độc lập với nhau. Gọi biến cố $A_i\colon$ \lq\lq Lần bắn thứ $i$ không trúng đích\rq\rq~ với $i\in \left\{1;2\right\}$. Trong các khẳng định nào sau đây, khẳng định nào đúng, khẳng định nào sai?
  \choiceTF
   {\True $A_1;A_2$ là hai biến cố độc lập}
  {Xác suất biến cố \lq\lq Cả hai lần bắn không trúng  đích\rq\rq~ là $0{,}5$}
  {\True Xác suất biến cố \lq\lq Lần bắn thứ nhất không trúng đích, lần bắn thứ hai trúng đích\rq\rq~ là $0{,}14$} 
  {\True Xác suất biến cố \lq\lq Có  ít nhất một lần bắn trúng đích\rq\rq~ là $0{,}94$}
 \loigiai{ 
  Ta có biến cố $A_i\colon$ \lq\lq Lần bắn thứ $i$  không trúng đích\rq\rq~ với $i=1,2$.\\
  Biến cố $\overline{A}_i\colon$ \lq\lq Lần  bắn thứ $i$ trúng đích\rq\rq~ với $i=1,2$.\\
  Theo giả thiết ta có $\mathrm{P}\left(A_1  \right)=0{,}2;\,\mathrm{P}\left(A_2 \right)=0{,}3;\mathrm{P}\left(\overline{A}_1\right)=0{,}8;\mathrm{P}\left(\overline{A}_2 \right)=0{,}7$.
  \begin{itemchoice}
   \itemch Vì các kết quả bắn là độc lập với nhau nên $A_1;\,A_2$ là hai biến cố  độc lập.
   \itemch Gọi biến cố $A\colon$ \lq\lq Cả hai lần bắn  đều không trúng đích\rq\rq.\\
   Ta có $A=A_1A_2$ và  $A_1;A_2$ là hai biến cố độc lập $\Rightarrow \mathrm{P}\left(A \right)=\mathrm{P}\left(A_1 \right)\cdot \mathrm{P}\left(A_2 \right)=0{,}2\cdot 0{,}3=0{,}06$.
   \itemch Gọi biến cố $B\colon$ \lq\lq Lần bắn thứ nhất không trúng đích,  lần bắn thứ hai trúng đích\rq\rq.\\
   Ta có $B=A_1\overline{A}_2$ và $A_1;\overline{A}_2$ là hai biến cố độc  lập $\Rightarrow \mathrm{P}\left(B \right)=\mathrm{P}\left(A_1 \right)\cdot\mathrm{P}\left(\overline{A}_2 \right)=0{,}2\cdot 0{,}7=0{,}14$.
   \itemch Gọi biến cố $C\colon$  \lq\lq Có ít nhất một lần bắn trúng đích\rq\rq.\\
   Khi đó, biến cố $\overline{C}\colon$ \lq\lq Cả hai lần bắn đều  không trúng đích\rq\rq.\\
   $\Rightarrow  \overline{C}=A \Rightarrow \mathrm{P}\left(\overline{C} \right)=\mathrm{P}\left(A \right)=0{,}06 \Rightarrow \mathrm{P}\left(C \right)=1-\mathrm{P}\left(\overline{C} \right)=0{,}94$.
  \end{itemchoice}
  }
\end{ex}

% nguồn: Ok_dungsai_lop11_Lan1
\begin{ex}%[1D9H1-3]%[ma:24d8e3a1]
 Một nhóm có $50$ người được phỏng vấn họ đã mua cành đào hay cây quất vào dịp Tết vừa qua, trong đó có $31$ người mua cành đào, $12$ người mua cây quất và $5$ người mua cả cành đào  và cây quất. Chọn ngẫu nhiên một người. Các khẳng định sau đây, khẳng định nào là đúng, khẳng định nào là sai?
  \choiceTF
   {Biến cố người đó mua đào và biến cố người đó mua quất là hai biến cố xung khắc}
  {\True  Xác suất để người đó mua cành đào là $\dfrac{31}{50}$, xác suất để người đó mua cành quất là $\dfrac{12}{50}$}
  {Xác suất để người đó mua cành đào hoặc mua cây quất là $\dfrac{19}{50}$} 
  {\True Xác suất  để người đó không mua cành đào và không mua cây quất là $\dfrac{6}{25}$}
  \loigiai{
  Gọi $A$ là biến cố \lq\lq Người đó mua cành đào\rq\rq, $B$ là biến cố \lq\lq Người đó  mua cây quất\rq\rq.
  \begin{itemchoice}
   \itemch Vì có $5$ người mua cả cành đào và cây quất nên biến cố $A$ và biến cố  $B$ không là hai biến cố xung khắc.
    \itemch Theo giả thiết $\mathrm{P}(A)=\dfrac{31}{50};\mathrm{P}(B)=\dfrac{12}{50};\mathrm{P}(AB)=\dfrac{5}{50}=\dfrac{1}{10}$.
   \itemch Lúc đó $A\cup B$ là biến cố \lq\lq Người đó mua cành đào hoặc mua cây  quất\rq\rq.\\
   Áp dụng công thức cộng ta  có $\mathrm{P}(A\cup B)=\mathrm{P}(A)+\mathrm{P}(B)-\mathrm{P}(AB)=\dfrac{31}{50}+\dfrac{12}{50}-\dfrac{5}{50}=\dfrac{38}{50}=\dfrac{19}{25}$.
   \itemch Lúc đó  $\overline{AB}$ là biến cố \lq\lq Người đó không mua cành đào và không mua cây quất\rq\rq.\\
   Vì $A\cup B$ biến cố đối của biến cố $\overline{AB}$  nên $\mathrm{P}(\overline{A}\,\overline{B})=1-\mathrm{P}(A\cup B)=1-\dfrac{38}{50}=\dfrac{12}{50}=\dfrac{6}{25}$.
  \end{itemchoice}
  }
\end{ex}

% nguồn: Ok_dungsai_lop11_Lan1
\begin{ex}%[1D9H1-3]%[ma:2be40dc5]
 Một hộp có chứa $5$ quả cầu trắng và $6$ quả cầu đen  cùng kích thước và khối lượng. Lấy ra ngẫu nhiên cùng một lúc $4$ quả cầu.
  \choiceTF
  {\True Xác suất để có hai  quả cầu trắng bằng $\dfrac{5}{11}$}
  {\True Xác suất để  có ít nhất $3$ quả cầu đen bằng $\dfrac{23}{66}$}
  {\True  Xác suất để cả $4$ quả cầu trắng bằng $\dfrac{1}{66}$}
  {Xác suất để không  có cầu trắng bằng $\dfrac{65}{66}$}
 \loigiai{ 
  Gọi $A$ biến cố \lq\lq Lấy ra ngẫu nhiên $4$ quả cầu trong đó  có $2$ quả cầu trắng\rq\rq.\\
  Gọi $B$ biến cố \lq\lq Lấy ra ngẫu nhiên $4$ quả cầu trong đó có ít nhất 3 quả  cầu đen\rq\rq.\\
  Gọi $C$ biến cố \lq\lq Lấy ra  ngẫu nhiên $4$ quả cầu toàn cầu trắng\rq\rq.\\
  Gọi $D$ biến  cố \lq\lq Lấy ra ngẫu nhiên $4$ quả trong đó không có quả cầu trắng\rq\rq.
  \begin{itemchoice}
   \itemch  \textbf{Đúng.}\\
   Ta có  $\mathrm{P}(A)=\dfrac{\mathrm{C}_5^2 \times \mathrm{C}_6^2}{\mathrm{C}_{11}^4}=\dfrac{5}{11}$.
   \itemch  \textbf{Đúng.}\\
   Ta có $\mathrm{P}(B)=\dfrac{\mathrm{C}_5^1 \times \mathrm{C}_6^3+\mathrm{C}_6^4}{\mathrm{C}_{11}^4}=\dfrac{23}{66}$. 
   \itemch  \textbf{Đúng.}\\
   Ta  có $\mathrm{P}(C)=\dfrac{\mathrm{C}_5^4}{\mathrm{C}_{11}^4}=\dfrac{1}{66}$.
   \itemch \textbf{Sai.}\\ 
    Ta có $\mathrm{P}(D)=\dfrac{\mathrm{C}_6^4}{\mathrm{C}_{11}^4}=\dfrac{1}{22}$.
  \end{itemchoice}
  }
\end{ex}

% nguồn: Ok_dungsai_lop11_Lan1
\begin{ex}%[1D9H1-3]%[ma:851d0089]
 Có hai hộp cùng chứa các quả cầu. Hộp thứ nhất có $7$ quả cầu đỏ, $5$ quả cầu xanh. Hộp  thứ $2$ có $6$ quả cầu đỏ, $4$ quả cầu xanh. Từ mỗi hộp lấy ra ngẫu nhiên $1$ quả cầu. Các khẳng định sau đây, khẳng định nào đúng, khẳng định nào sai?
  \choiceTF
  {\True Xác suất để quả cầu lấy ra từ hộp thứ nhất có màu đỏ là  $\dfrac{7}{12}$}
  {Xác suất để hai quả  cầu lấy ra cùng màu đỏ là $\dfrac{1}{2}$}
  {Xác suất để 2quả cầu lấy ra có ít nhất $1$ quả màu  đỏ là $\dfrac{3}{8}$}
  {\True  Xác suất để $2$ quả cầu lấy ra cùng màu là $\dfrac{31}{60}$}
 \loigiai{ 
  \begin{itemchoice}
   \itemch Hộp thứ  nhất có tất cả $12$ quả cầu, trong đó có $7$ quả cầu đỏ.\\
   Gọi $A$  là biến cố \lq\lq Quả cầu lấy ra từ hộp thứ nhất màu đỏ\rq\rq~ ta có $\mathrm{P}(A)=\dfrac{7}{12}$.
   \itemch Gọi $B$  là biến cố \lq\lq Quả cầu lấy ra từ hộp thứ hai màu đỏ\rq\rq~ ta có $\mathrm{P}(B)=\dfrac{6}{10}=\dfrac{3}{5}$.\\
   Gọi $C$ là biến cố \lq\lq Hai quả cầu lấy ra cùng màu đỏ\rq\rq~ ta có $C=AB$.\\ 
    Vì hai biến cố $A, B$ độc lập nên $\mathrm{P}(C)=\mathrm{P}(AB)=\mathrm{P}(A)\cdot\mathrm{P}(B)=\dfrac{7}{12}\cdot\dfrac{3}{5}=\dfrac{7}{20}$.
   \itemch Gọi $D$ là biến cố \lq\lq Hai quả cầu lấy ra có ít  nhất $1$ quả màu đỏ\rq\rq~ ta có $D=A\cup B$.\\
   Theo công thức cộng xác suất  ta có $\mathrm{P}(D)=\mathrm{P}(A\cup B)=\mathrm{P}(A)+\mathrm{P}(B)-\mathrm{P}(AB)=\dfrac{7}{12}+\dfrac{3}{5}-\dfrac{7}{20}=\dfrac{5}{6}$.
   \itemch Gọi $F$ là biến cố \lq\lq Hai quả cầu lấy ra có cùng màu  xanh\rq\rq. Ta có $F=\overline{A}\,\overline{B}$.\\
   Vi $A,B$ độc lập nên  $\overline{A},\overline{B}$ độc lập. Do đó
   $$\mathrm{P}(F)=\mathrm{P}(\overline{A}\,\overline{B})=\mathrm{P}(\overline{A})\cdot\mathrm{P}(\overline{B})=\left[1-\mathrm{P}(A) \right]\cdot\left[1-\mathrm{P}(B) \right]=\left[1-\dfrac{7}{12}  \right]\cdot\left[1-\dfrac{3}{5} \right]=\dfrac{1}{6}.$$
   Gọi $X$ là biến  cố \lq\lq Hai quả cầu lấy ra cùng màu\rq\rq~ ta thấy $X=C\cup F;C\cap F=\varnothing$ nên
    $$\mathrm{P}(X)=\mathrm{P}(C)+\mathrm{P}(F)=\dfrac{7}{20}+\dfrac{1}{6}=\dfrac{31}{60}.$$
  \end{itemchoice}
 } 
\end{ex}

% nguồn: Ok_dungsai_lop11_Lan1
\begin{ex}%[1D9H1-3]%[ma:d2d28f61]
 Một chiếc máy có hai động cơ I và II hoạt động độc lập với nhau. Xác suất để động cơ I và động cơ II chạy tốt lần  lượt là $0{,}8$ và $0{,}9$. Các khẳng định sau đây, khẳng định nào đúng, khẳng định nào sai?
 \choiceTF 
  {\True Xác suất để cả hai động cơ đều chạy  tốt là $0{,}72$}
  {Xác suất  để cả hai động cơ đều không chạy tốt là $0{,}63$}
  {Xác suất để có ít nhất một động cơ  chạy tốt là $0{,}94$}
   {\True Xác suất để một trong hai động chay tốt là $0{,}26$}
 \loigiai{ 
  Gọi $A$ là biến cố \lq\lq Động cơ I chạy tốt\rq\rq;\\ 
  $B$ là  biến cố \lq\lq Động cơ II chạy tốt\rq\rq;\\
  $C$ là biến cố \lq\lq Cả hai động cơ đều  chạy tốt\rq\rq;\\
  $D$ là biến cố \lq\lq Cả  hai động cơ đều không chạy tốt\rq\rq;\\
  $E$ là biến cố \lq\lq Có ít nhất một  động cơ chạy tốt\rq\rq;\\
  $F$ là biến cố \lq\lq Chỉ một trong hai động  cơ chạy tốt\rq\rq.
  \begin{itemchoice}
   \itemch Ta thấy $A,\,B$ là hai biến cố độc lập  với nhau và $C=A\cap B$.
   Áp dụng quy tắc nhân xác suất  ta có $\mathrm{P}(C)=\mathrm{P}(AB)=\mathrm{P}(A)\cdot\mathrm{P}(B)=0{,}8\cdot 0{,}9=0{,}72$.
    \itemch Ta có $\mathrm{P}(\overline{A})=1-0{,}8=0{,}2$ và $\mathrm{P}(\overline{B})=1-0{,}9=0{,}1$.\\
   Vì $A,\,B$ là hai biến cố độc  lập nên $\overline{A},\,\overline{B}$ cũng độc lập và $D=\overline{A}\cap\overline{B}$.\\
   Do đó  $\mathrm{P}(D)=\mathrm{P}(\overline{A})\cdot\mathrm{P}(\overline{B})=0{,}2\cdot 0{,}1=0{,}02$.
   \itemch Gọi $E$ là biến cố  \lq\lq Có ít nhất một động cơ chạy tốt\rq\rq.\\
   Biến cố đối của $E$ là \lq\lq Cả hai  động cơ chạy không tốt\rq\rq~ $\Rightarrow \overline{E}=D=\overline{A}\cap\overline{B}$.\\
   Do đó  $\mathrm{P}(E)=1-\mathrm{P}(D)=1-\mathrm{P}(\overline{A}\,\overline{B})=1-\mathrm{P}(\overline{A})\cdot\mathrm{P}(\overline{B})=1-0{,}2\cdot 0{,}1=0{,}98$.
   \itemch Ta có  $F=\left(A\cap \overline{B}\right) \cup \left(\overline{A}\cap B\right)$.\\
   Ta có $A\cap \overline{B}$ và $\overline{A}\cap B$ xung khắc và $\overline{A},\,B,\,A,\,\overline{B}$ độc lập  nên ta có
    \allowdisplaybreaks
   \begin{eqnarray*}
    \mathrm{P}(F) &=&  \mathrm{P}\left((A\cap \overline{B})\cup (\overline{A}\cap B) \right)=\mathrm{P}(A\cap \overline{B})+\mathrm{P}(\overline{A}\cap B) \\
    &=& \mathrm{P}(A)\cdot \mathrm{P}(\overline{B})+\mathrm{P}(\overline{A})\cdot\mathrm{P}(B)=0{,}8\cdot  0{,}1+0{,}2\cdot 0{,}9=0{,}26.
   \end{eqnarray*}
  \end{itemchoice}
 } 
\end{ex}

% nguồn: Ok_dungsai_lop11_Lan1
\begin{ex}%[1D9H1-3]%[ma:3858b0c7]
 Một hộp có chứa $6$ bút  mực xanh và $4$ bút mực đỏ cùng loại, cùng kích thước và khối lượng. Lấy ra ngẫu nhiên đồng thời $3$ bút từ hộp. Gọi $A$ là biến cố \lq\lq Ba bút lấy ra đều là bút mực xanh\rq\rq. $B$ là biến cố \lq\lq Ba bút lấy ra đều là bút mực đỏ\rq\rq. Khi đó:
  \choiceTF
  {Có $30$ kết quả thuận lợi cho biến cố $A$} 
  {\True  Có $4$ kết quả thuận lợi cho biến cố $B$}
  {\True  Xác suất của biến cố $A$ bằng $\dfrac{1}{6}$}
  {\True Xác suất  của biến cố $B$ bằng $\dfrac{1}{30}$}
 \loigiai{ 
  \begin{itemchoice}
    \itemch \textbf{Sai.}\\
   Số kết  quả thuận lợi cho biến cố $A$ là $\mathrm{C}_6^3=20$.
   \itemch \textbf{Đúng.}\\ 
   Số  kết quả thuận lợi cho biến cố $B$ là $\mathrm{C}_4^3=4$.
    \itemch \textbf{Đúng.}\\
   Xác suất  của biến cố $A$ là $\mathrm{P}(A)=\dfrac{\mathrm{C}_6^3}{\mathrm{C}_{10}^3}=\dfrac{1}{6}$.
   \itemch \textbf{Đúng.}\\ 
   Xác suất của biến  cố $B$ là $\mathrm{P}(B)=\dfrac{\mathrm{C}_4^3}{\mathrm{C}_{10}^3}=\dfrac{1}{30}$.
  \end{itemchoice}
 } 
\end{ex}

% nguồn: Ok_dungsai_lop11_Lan1
\begin{ex}%[1D9H1-3]%[ma:fb95c5a7]
 Cho $A$ và $B$ là hai  biến cố độc lập với nhau, biết $\mathrm{P}(A)=0{,}2$, $\mathrm{P}(B)=0{,}3$. Khi đó
 \choiceTF 
  {\True $\mathrm{P}(AB)=0{,}06$} 
  {$\mathrm{P}(A\overline{B})=0{,}12$} 
   {\True $\mathrm{P}(\overline{A}\overline{B})=0{,}56$}
  {\True $\mathrm{P}(\overline{A}B)=0{,}24$} 
 \loigiai{ 
  \begin{itemchoice}
   \itemch \textbf{Đúng.}\\ 
   Vì $A$ và $B$ là hai  biến cố độc lập nên $\mathrm{P}(AB)=\mathrm{P}(A)\cdot \mathrm{P}(B)=0{,}2\cdot 0{,}3=0{,}06$.
    \itemch \textbf{Sai.}\\
   Vì $A$ và $B$ là hai biến cố độc lập nên $A$,  $\overline{B}$ độc lập.\\
    Suy ra $\mathrm{P}(A\overline{B})=\mathrm{P}(A)\cdot \mathrm{P}(\overline{B})=0{,}2\cdot 0{,}7=0{,}14$.
   \itemch \textbf{Đúng.}\\ 
   Vì $A$ và $B$ là hai biến cố độc lập nên $\overline{A}$,  $\overline{B}$ độc lập.\\
   Suy ra $\mathrm{P}(\overline{A}\overline{B})=\mathrm{P}(\overline{A})\cdot \mathrm{P}(\overline{B})=0{,}8\cdot  0{,}7=0{,}56$.
   \itemch \textbf{Đúng.}\\ 
   Vì $A$  và $B$ là hai biến cố độc lập nên $\overline{A}$, $B$ độc lập.\\
   Suy ra  $\mathrm{P}(\overline{A}B)=\mathrm{P}(\overline{A})\cdot \mathrm{P}(B)=0{,}8.0{,}3=0{,}24$.
  \end{itemchoice}
 } 
\end{ex}

% nguồn: Ok_dungsai_lop11_Lan1
\begin{ex}%[1D9H1-3]%[ma:4e320cc9]
 Một người vừa gieo một con xúc xắc để ghi lại số  chấm xuất hiện, sau đó người này tiếp tục chọn ngẫu nhiên một lá bài từ bộ bài $52$ lá.
  \choiceTF
  {\True Xác suất để số chấm của con  xúc xắc lớn nhất bằng $\dfrac{1}{6}$}
  {\True Xác suất để chọn được một lá bài tây (quân  bài có hình người) bằng $\dfrac{3}{13}$}
  {\True Xác suất để số chấm trên con xúc xắc là lớn nhất và chọn được một lá bài  tây bằng $\dfrac{1}{26}$}
  {Xác suất để số  chấm trên con xúc xắc và số của lá bài là giống nhau bằng $\dfrac{1}{16}$}
  \loigiai{
  Gọi $A$ là biến cố  \lq\lq Số chấm của xúc xắc lớn nhất\rq\rq.\\
   Gọi $B$ là biến cố: \lq\lq Chọn được một lá bài tây\rq\rq.
  \begin{itemchoice}
   \itemch \textbf{Đúng.}\\ 
   Ta  có $A=\{6\} \Rightarrow n(A)=1\Rightarrow \mathrm{P}(A)=\dfrac{1}{6}$.
   \itemch \textbf{Đúng.}\\ 
   Ta biết bộ bài $52$ lá thì có $12$ lá bài tây, nên xác suất chọn được một  lá bài tây là $\mathrm{P}(B)=\dfrac{3}{13}$.
    \itemch \textbf{Đúng.}\\
   Dễ thấy  $A$, $B$ là hai biến cố độc lập.\\
   Suy ra $\mathrm{P}(AB)=\mathrm{P}(A)\cdot  \mathrm{P}(B)=\dfrac{1}{6} \cdot \dfrac{3}{13}=\dfrac{1}{26}$.
   \itemch \textbf{Sai.}\\ 
   Để thu được số chấm trên con xúc xắc và số của lá  bài giống nhau thì ta có $6$ cách để có được số chấm một con xúc xắc, ứng với mỗi cách đó thì có đúng $4$ cách tìm được lá bài thoả mãn.\\
   Việc gieo xúc xắc và rút ngẫu nhiên lá bài là độc lập.\\ 
   Gọi $X$ là biến cố cần tính xác suất, ta có  $\mathrm{P}(X)=\dfrac{6}{6} \cdot \dfrac{4}{52}=\dfrac{1}{13}$.
  \end{itemchoice}
 } 
\end{ex}

% nguồn: Ok_dungsai_lop11_Lan1
\begin{ex}%[1D9V1-3]%[ma:3e27f2de]
 Trong kì thi THPT Quốc Gia năm $2023$ có môn thi bắt buộc  là môn Toán. Môn thi này thi dưới hình thức trắc nghiệm với bốn phương án trả lời $A,\,B,\,C,\,D$. Mỗi câu trả lời đúng được cộng $0{,}2$ điểm; mỗi câu trả lời sai bị trừ $0{,}1$ điểm. Bạn Việt vì học rất kém môn Toán nên chọn ngẫu nhiên cả $50$ câu trả lời. Các khẳng định sau đây khẳng định nào đúng, khẳng định nào sai?
 \choiceTF 
  {\True Xác  suất để bạn Việt làm đúng ở một câu là $\dfrac{1}{4}$}
  {Xác suất để bạn Việt làm sai  ở một câu là $\dfrac{1}{3}$}
  {Bạn Việt làm được $7$ điểm, khi  đó bạn Việt làm đúng $30$ câu và sai $20$ câu}
  {Xác suất để bạn Việt đạt được $4$ điểm môn Toán trong  kì thi trên là $2{,}5\cdot10^{-6}$}
  \loigiai{
  \begin{itemchoice}
   \itemch Trong một câu có bốn phương án trả lời, có đúng một  phương án đúng.\\
   Suy ra xác suất  để bạn Việt làm đúng ở một câu là $\dfrac{1}{4}$.
   \itemch Trong  một câu có bốn phương án trả lời, có ba phương án sai.\\
   Suy ra xác suất để bạn Việt làm sai  ở một câu là $\dfrac{3}{4}$.
   \itemch Gọi $x$ là số câu đúng Việt chọn để được $7$ điểm, ta có  $0{,}2\cdot x-(50-x)\cdot 0{,}1=7 \Leftrightarrow x=40$.\\
    Vậy bạn Việt làm được $7$ điểm, khi đó bạn Việt làm đúng $40$ câu và sai $10$ câu.
   \itemch Gọi $x$ là số câu đúng Việt chọn để được $4$ điểm, ta có  $0{,}2\cdot x-(50-x)\cdot 0{,}1=4 \Leftrightarrow x=30$.\\
   Vậy xác suất  Việt đạt $4$ điểm môn Toán trong kì thi trên là
   $$\mathrm{P}=\mathrm{C}_{50}^{30}\cdot\left(\dfrac{1}{4} \right)^{30}\cdot\left(\dfrac{3}{4} \right)^{20}=1{,}3\cdot  10^{-7}.$$
  \end{itemchoice}
  }
\end{ex}

% nguồn: Ok_dungsai_lop11_Lan1
\begin{ex}%[1D9V1-3]%[ma:abe30036]
 Trong  đề kiểm tra $15$ phút môn Toán có $20$ câu trắc nghiệm. Mỗi câu trắc nghiệm có $4$ phương án trả lời, trong đó chỉ có một phương án trả lời đúng. Bình giải chắc chắn đúng $10$ câu, $10$ câu còn lại lựa chọn ngẫu nhiên đáp án. Biết rằng mỗi câu trả lời đúng được $0{,}5$ điểm, trả lời sai không bị trừ điểm. Các khẳng định sau đây, khẳng định nào đúng, khẳng định nào sai?
 \choiceTF 
  {\True Bình chắc chắn được $5$  điểm}
  {\True Xác  suất trả lời sai ở một câu là $\dfrac{3}{4}$}
  {Xác suất để Bình đạt  đúng $8$ điểm là $\left(\dfrac{1}{4} \right)^6\cdot\left(\dfrac{3}{4} \right)^4$}
  {Xác suất để Bình đạt  được từ $9$ điểm trở lên nhỏ hơn 0{,}0004}
  \loigiai{
  \begin{itemchoice}
   \itemch Bình giải chắc chắn đúng $10$ câu nên Bình được chắc chắn $5$  điểm.
   \itemch Xác suất khi đánh ngẫu nhiên sai một câu  trắc nghiệm là $\dfrac{3}{4}$.
   \itemch Để Bình đạt được  đúng $8$ điểm thì trong $10$ câu còn lại lựa chọn ngẫu nhiên đáp án phải đúng $6$ câu, sai $4$ câu.\\
   Xác suất khi đánh ngẫu nhiên đúng  một câu trắc nghiệm là $\dfrac{1}{4}$.\\
   Xác suất khi đánh ngẫu nhiên sai một câu  trắc nghiệm là $\dfrac{3}{4}$.\\
   Chọn  $6$ câu trắc nghiệm để đáp đúng từ $10$ câu trắc nghiệm có $\mathrm{C}_{10}^6$.\\
   Vậy xác suất để Bình đạt được đúng $8$ điểm là $\mathrm{C}_{10}^6\cdot\left(\dfrac{1}{4} \right)^6\cdot\left(\dfrac{3}{4}  \right)^4$.
   \itemch Để Bình  đạt được từ $9$ điểm trở lên thì trong $10$ câu còn lại lựa chọn ngẫu nhiên đáp án phải đúng ít nhất $8$ câu.\\
   Vậy xác suất để Bình đạt được từ $9$ điểm trở lên  là
   $$\mathrm{C}_{10}^8\cdot\left(\dfrac{1}{4} \right)^8\cdot\left(\dfrac{3}{4} \right)^2+\mathrm{C}_{10}^9\cdot\left(\dfrac{1}{4} \right)^9\cdot\left(\dfrac{3}{4}  \right)^1+\left(\dfrac{1}{4} \right)^{10}=\dfrac{109}{262144}>0{,}0004.$$
  \end{itemchoice}
  }
\end{ex}

% nguồn: Ok_dungsai_lop11_Lan1
\begin{ex}%[1D9V1-3]%[ma:9a0efb9c]
 Trên một bảng quảng cáo, người ta mắc hai hệ thống bóng đèn. Hệ thống I gồm $2$ bóng mắc nối tiếp, hệ thống II gồm $2$ bóng mắc song song. Khả năng bị hỏng của mỗi bóng đèn sau $6$ giờ thắp  sáng liên tục là $0{,}15$. Biết tình trạng của mỗi bóng đèn là độc lập.
  \choiceTF
  {\True Xác suất để hệ  thống II bị hỏng (không sáng) bằng $0{,}0225$}
  {\True Xác suất để hệ  thống II hoạt động bình thường bằng $0{,}9775$}
  {Xác suất để  hệ thống I bị hỏng (không sáng) bằng $0{,}5775$}
  {Xác suất để  cả hai hệ thống bị hỏng (không sáng) (kết quả được làm tròn đến hàng phần trăm nghìn) bằng $\approx 0{,}02624$}
 \loigiai{ 
  \begin{itemchoice}
    \itemch \textbf{Đúng.}\\
   Hệ thống  II gồm $2$ bóng được mắc song song nên nó chỉ hỏng khi cả hai bóng đều hỏng.\\
    Gọi $B$ là biến cố \lq\lq Hệ thống II bị hỏng\rq\rq, ta có $\mathrm{P}(B)=0{,}15\cdot 0{,}15=0{,}0225$.
   \itemch  \textbf{Đúng.}\\
    Xác suất để hệ thống II hoạt động bình thường là $\mathrm{P}\overline{B})=1-0{,}0225=0{,}9775$.
    \itemch \textbf{Sai.}\\
   Hệ thống I chỉ  hoạt động bình thường khi cả hai bóng bình thường.\\
   Gọi $A$ là biến cố \lq\lq Hệ thống $I$ bị  hỏng\rq\rq. Khi đó xác suất để hệ thống $I$ hoạt động bình thường là $\mathrm{P}\overline{A})=0{,}85\cdot 0{,}85=0{,}7225$.\\
   Suy ra  $\mathrm{P}(A)=1-0{,}7225=0{,}2775$
    \itemch \textbf{Sai.}\\
   Xác suất để cả hai hệ thống  I, II đều bị hỏng là
   \[ 
   \mathrm{P}(AB)=\mathrm{P}(A)\cdot \mathrm{P}(B)=0{,}2775\cdot  0{,}0225=\dfrac{999}{160000} \approx 0{,}00624.
   \] 
  \end{itemchoice}
 } 
\end{ex}

% nguồn: Ok_dungsai_lop11_Lan1
\begin{ex}%[1D9V1-3]%[ma:8ece1693]
 Một người vừa gieo một con xúc xắc để ghi lại số chấm xuất hiện, sau đó người này tiếp tục chọn ngẫu nhiên một lá bài từ bộ  bài $52$ lá.
 \choiceTF 
  {\True Gọi $A$  là biến cố: \lq\lq Số chấm của xúc xắc lớn nhất\rq\rq, khi đó: $\mathrm{P}(A)=\dfrac{1}{6}$}
  {\True Gọi $B$ là biến cố: \lq\lq  Chọn được một lá bài tây\rq\rq, khi đó: $\mathrm{P}(B)=\dfrac{3}{13}$}
  {\True Xác suất để số chấm trên con xúc xắc là lớn nhất và chọn  được một lá bài tây bằng $\dfrac{1}{26}$}
  {Xác suất để số chấm trên con xúc xắc và số của lá bài là giống nhau bằng $\dfrac{1}{16}$} 
  \loigiai{
  Gọi $A$ là biến  cố: \lq\lq Số chấm của xúc xắc lớn nhất\rq\rq, $B$ là biến cố: \lq\lq Chọn được một lá bài tây\rq\rq. Dễ thấy $A$, $B$ là hai biến cố độc lập.
  \begin{itemchoice}
   \itemch Đúng. Ta  có: $A=\{6\}\Rightarrow n(A)=1\Rightarrow\mathrm{P}(A)=\dfrac{1}{6}$.
   \itemch Đúng. Ta  biết bộ bài 52 lá thì có 12 lá bài tây, nên xác suất chọn được một lá bài tây là $\mathrm{P}(B)=\dfrac{3}{13}$.
    \itemch Đúng. Suy ra $\mathrm{P}(AB)=\mathrm{P}(A)\cdot\mathrm{P}(B)=\dfrac{1}{6}\cdot\dfrac{3}{13}=\dfrac{1}{26}$.
   \itemch  Sai.
   Để thu được số chấm trên con xúc xắc và số của lá bài giống nhau thì ta có 6 cách để có được số chấm một con xúc xắc, ứng với mỗi cách đó thì có đúng 4 cách tìm được lá bài  thoả mãn.
   Việc gieo xúc xắc và rút ngẫu nhiên lá bài là độc lập.\\ 
    Gọi $X$ là biến cố cần tính xác suất, ta có: $\mathrm{P}(X)=\dfrac{6}{6}\cdot\dfrac{4}{52}=\dfrac{1}{13}$.
 \end{itemchoice}}
\end{ex}


%%%%%%%%%%%%%%%% TRẢ LỜI NGẮN %%%%%%%%%%%%%%%%

% nguồn: Ok_TraLoiNgan_Lop11_Lan1
\begin{ex}%[1D9H1-3]%[ma:3fbf0ea0]
	Ba xạ thủ lần lượt bắn vào  một bia. Xác suất để xạ thủ thứ nhất, thứ hai, thứ ba bắn trúng đích lần lượt là $0{,}8$; $ 0{,}6$; $0{,}5$. Tính xác suất để có đúng hai người bắn trúng đích.
	\shortans{0,46} 
	\loigiai{ 
		Gọi $A_i(1\leq i\leq 3, i\in\mathbb{N})$ lần lượt là biến cố \lq\lq Xạ thủ thứ $i$  bắn trúng đích\rq\rq.\\
		Xác suất để xạ  thủ thứ nhất, thứ hai, thứ ba bắn trúng đích lần lượt là $\mathrm{P}(A_1)=0{,}8$; $\mathrm{P}(A_2)=0{,}6$; $P(A_3)=0{,}5$.\\
		Gọi $X$ là biến cố \lq\lq Có  đúng hai xạ thủ bắn trúng đích\rq\rq.\\
		Ta  có
		\begin{eqnarray*}
			 \mathrm{P}(X)
			&=& \mathrm{P}(A_1)\cdot \mathrm{P}(A_2)\cdot \mathrm{P}(\overline{A_3})+\mathrm{P}(A_1)\cdot \mathrm{P}(\overline{A_2})\cdot \mathrm{P}(A_3)+\mathrm{P}(\overline{A_1})\cdot \mathrm{P}(A_2)\cdot \mathrm{P}(A_3)\\ 
			&=&  0{,}8\cdot 0{,}6\cdot 0{,}5+0{,}8\cdot 0{,}4\cdot 0{,}5+0{,}2\cdot 0{,}6\cdot 0{,}5\\
			&=& 0{,}46. 
		\end{eqnarray*}
	} 
\end{ex}

% nguồn: Ok_TraLoiNgan_Lop11_Lan1
\begin{ex}%[1D9H1-3]%[ma:9e59afa5]
	Một bệnh truyền nhiễm có xác suất lây bệnh là $0{,}9$ nếu tiếp xúc với người bệnh mà không đeo khẩu trang; là $0{,}15$ nếu tiếp xúc với người bệnh mà có đeo khẩu trang. Anh Hà tiếp xúc với một người bệnh hai lần, trong đó có một lần đeo khẩu trang và một lần không đeo khẩu trang. Khả năng anh Hà bị lây bệnh từ người bệnh mà anh tiếp  xúc đó bằng bao nhiêu phần trăm?
	\shortans{13,5} 
	 \loigiai{
		Gọi $A$ là biến cố \lq\lq Anh Hà bị lây bệnh từ người bệnh nếu tiếp xúc với  người bệnh mà không đeo khẩu trang\rq\rq.\\
		Suy  ra $\mathrm{P}(A)=0{,}9$.\\
		Gọi $B$ là biến cố \lq\lq Anh Hà bị lây bệnh từ  người bệnh nếu tiếp xúc với người bệnh mà có đeo khẩu trang\rq\rq.\\
		 Suy ra $\mathrm{P}(B)=0{,}15$.\\
		Vì $A$ và $B$ là 2 biến cố độc lập nên xác suất của biến cố "anh Hà bị lây bệnh  từ người bệnh mà anh tiếp xúc đó'' là
		\[ 
		 \mathrm{P}(AB)=0{,}9\times 0{,}15=0{,}135.
		\] 
	} 
\end{ex}

% nguồn: Ok_TraLoiNgan_Lop11_Lan1
\begin{ex}%[1D9H1-3]%[ma:8ccb9632]
	Jenny, Jim đang trò đang trò chuyện với  Merry về việc có nên đi dự tiệc hay không. Xác suất Jenny sẽ tham dự là $0{,}4$ và xác suất Jim sẽ tham dự là $0{,}6$. Nhưng hôm đó Merry có việc bận nên khả năng không tham dự bữa tiệc là $0{,}8$ . Tính xác suất để ba người bạn cùng tham dự. (Làm tròn kết quả đến hàng phần trăm).
	\shortans{0,05} 
	\loigiai{ 
		Gọi $A$ là biến cố \lq\lq Jenny tham gia  bữa tiệc\rq\rq.\\
		Gọi $B$ là biến cố \lq\lq Merry tham gia  bữa tiệc\rq\rq. \\
		Gọi $C$  là biến cố \lq\lq Jim tham gia bữa tiệc\rq\rq.\\
		Các biến cố $A$, $B$,  $C$ đôi một độc lập với nhau.\\
		Xác suất để  ba người bạn cùng tham dự là
		\[ 
		 \mathrm{P}(ABC) = \mathrm{P}(A) \cdot \mathrm{P}(B) \cdot \mathrm{P}(C) = 0{,}4 \cdot 0{,}6 \cdot (1- 0{,}8) = 0, 048 \approx 0{,}05.
		 \]
	 }
\end{ex}

% nguồn: Ok_TraLoiNgan_Lop11_Lan1
\begin{ex}%[1D9H1-3]%[ma:41899a25]
	Trong một trận đấu bóng đá quan trọng ở vòng đấu loại trực tiếp, khi trận đấu buộc phải giải quyết bằng loạt sút luân lưu $11$ m, huấn luyện viên đội $X$ đưa danh sách lần lượt $5$ cầu thủ có xác suất sút luân lưu $11$ m thành công là $0{,}8$; $0{,}8$; $0{,}76$; $0{,}72$; $0{,}68$. Tìm xác suất để chỉ có cầu thủ cuối cùng sút trượt luân lưu (kết quả được làm tròn đến hàng phần  trăm).
	\shortans{0,11} 
	\loigiai{ 
		Gọi $A_i (1\le i\le 5,i\in \mathbb{N})$ là biến cố \lq\lq Cầu  thủ thứ $i$ của đội $X$ sút luân lưu thành công\rq\rq.
		Xác suất cần tìm  là
		\begin{eqnarray*}
			\mathrm{P}\left(A_1 A_2 A_3 A_4 \overline{A_5} \right) & = & \mathrm{P}\left(A_1\right) \cdot \mathrm{P}\left(A_2\right) \cdot \mathrm{P}\left(A_3\right) \cdot \mathrm{P}\left(A_4\right) \cdot \mathrm{P}\left(\overline{A_5}\right)  \\
			& = & 0{,}8 \cdot  0{,}8 \cdot 0{,}76 \cdot 0{,}72 \cdot (1-0{,}68)\\
			& \approx &  0{,}11.
		\end{eqnarray*}
	 }
\end{ex}

% nguồn: Ok_TraLoiNgan_Lop11_Lan1
\begin{ex}%[1D9H1-3]%[ma:620d6379]
	Một hộp có $10$ quả bóng bàn trong đó có $6$  quả mới. Người ta lấy ra ngẫu nhiên $5$ quả để thi đấu. Biết xác suất của biến cố lấy được ít nhất $2$ quả bóng mới bằng $\dfrac{a}{b}$, với $a$, $b$ là các số tự nhiên và phân số $\dfrac{a}{b}$ tối giản. Giá trị của $a + b$ bằng bao nhiêu?
	 \shortans{83}
	 \loigiai{
		Gọi  $A$ là biến cố \lq\lq Lấy ra ngẫu nhiên $5$ quả trong đó có ít nhất $2$ quả bóng mới\rq\rq.\\
		Xác suất  của biến cố $A$ là
		 \[
		\mathrm{P}(A)=\dfrac{\mathrm{C}_6^2 \times \mathrm{C}_4^3 + \mathrm{C}_6^3 \times \mathrm{C}_4^2  + \mathrm{C}_6^4 \times \mathrm{C}_4^1+\mathrm{C}_6^5}{\mathrm{C}_{10}^5}=\dfrac{41}{42}.
		\] 
		Suy ra $a = 41$, $b =42 \Rightarrow a + b =  83$.
	 }
\end{ex}

% nguồn: Ok_TraLoiNgan_Lop11_Lan1
\begin{ex}%[1D9H1-3]%[ma:8f13fb57]
	Một lô hàng có $20$ sản phẩm giống nhau trong đó có $4$ sản phẩm không đạt chất lượng còn lại là sản phẩm đạt chất lượng tốt. Mỗi lần kiểm tra, người ta chọn ra ngẫu nhiên $2$ sản phẩm. Tính xác suất để lấy ra được ít nhất một sản phẩm  tốt. (Làm tròn kết quả đến hàng phần mười).
	\shortans{96,8} 
	 \loigiai{
		Gọi $A$ là biến cố \lq\lq Lấy ra $2$ sản  phẩm được ít nhất một sản phẩm tốt\rq\rq.\\
		Biến cố $A$ xảy ra khi lấy ra được $1$ sản phẩm tốt và một sản phẩm không đạt chất lượng hoặc lấy  ra được $2$ sản phẩm đều tốt.\\
		 Xác suất của biến cố $A$ là $\mathrm{P}(A)=\dfrac{\mathrm{C}_{16}^1 \times \mathrm{C}_4^1 \times \mathrm{C}_{16}^2}{\mathrm{C}_{20}^2}=\dfrac{92}{95} = \approx 96{,}8$.
	 }
\end{ex}

% nguồn: Ok_TraLoiNgan_Lop11_Lan1
\begin{ex}%[1D9H1-3]%[ma:2ea353ea]
	Tung đồng thời một đồng xu và một con xúc xắc $12$ mặt (từ $1$ đến $12$). Xác suất để đồng xu xuất hiện mặt ngửa  và con xúc xắc xuất hiện mặt là bội của $3$ bằng bao nhiêu? (Làm tròn kết quả đến hàng phần trăm)
	 \shortans{0,17}
	 \loigiai{
		Gọi $A$ là biến cố \lq\lq Đồng xu xuất hiện mặt ngửa\rq\rq.  Suy ra $\mathrm{P}(A) = \dfrac{1}{2}$.\\
		Gọi $C$ là biến cố \lq\lq Con xúc xắc  xuất hiện mặt là bội của $3$\rq\rq. Suy ra $\mathrm{P}(C)=\dfrac{4}{12}=\dfrac{1}{3}$.\\
		Xác suất xuất  hiện mặt ngửa và mặt là bội của $3$ là
		 \[
		\mathrm{P}(AC)=\mathrm{P}(A)\cdot  \mathrm{P}(C)=\dfrac{1}{2} \cdot \dfrac{1}{3}=\dfrac{1}{6} \approx 0{,}17.
		 \]
	} 
\end{ex}

% nguồn: Ok_TraLoiNgan_Lop11_Lan1
\begin{ex}%[1D9H1-3]%[ma:e93612df]
	Hai xạ thủ cùng bắn vào bia một cách độc lập với nhau. Xác suất bắn trúng bia của xạ thủ thứ nhất bằng $\dfrac{1}{2}$, xác suất bắn trúng bia của xạ  thủ thứ hai bằng $\dfrac{1}{3}$. Xác suất để cả hai xạ thủ đều bắn không trúng bia bằng bao nhiêu? (Làm tròn kết quả đến hàng phần trăm).
	 \shortans{0,33}
	 \loigiai{
		Gọi $A$ là biến cố \lq\lq Xạ thủ thứ nhất bắn trúng bia\rq\rq.\\ 
		Gọi $B$ là  biến cố \lq\lq Xạ thủ thứ hai bắn trúng bia\rq\rq.\\
		Khi đó $A$, $B$ là  các biến cố độc lập.\\
		Ta có $\mathrm{P}(\overline{A})=1-\dfrac{1}{2}=\dfrac{1}{2}, \mathrm{P}(\overline{B})=1-\dfrac{1}{3}=\dfrac{2}{3}$.\\ 
		Xác suất của biến cố $\overline{A}\overline{B}$ là $\mathrm{P}(\overline{A}\overline{B})=\mathrm{P}(\overline{A})\cdot \mathrm{P}(\overline{B})=\dfrac{1}{2} \cdot \dfrac{2}{3}=\dfrac{1}{3} \approx  0{,}33$.
	 }
\end{ex}

% nguồn: Ok_TraLoiNgan_Lop11_Lan1
\begin{ex}%[1D9H1-3]%[ma:1ad7056a]
	Nhà trường muốn chọn một đội văn nghệ  có đủ cả nam và nữ gồm $12$ em đi biểu diễn từ một nhóm học sinh gồm $10$ nam sinh và $8$ nữ sinh. Tính xác xuất để đội văn nghệ được chọn có ít nhất 6 bạn nữ. (Làm tròn kết quả đến hàng phần mười).
	 \shortans{43,7}
	 \loigiai{
		Gọi $B$  là biến cố \lq\lq Chọn đội văn nghệ gồm 12 bạn có đủ cả nam và nữ trong đó có ít nhất 6 bạn nữ'\rq\rq.\\
		 Suy ra
		 \[
		\mathrm{P}(B)=\dfrac{\mathrm{C}_{10}^6 \times \mathrm{C}_8^6+\mathrm{C}_{10}^5 \times \mathrm{C}_8^7+\mathrm{C}_{10}^4 \times  \mathrm{C}_8^8}{\mathrm{C}_{18}^{12}}=\dfrac{193}{442} = \approx 43{,}7.
		\] 
	 }
\end{ex}

% nguồn: Ok_TraLoiNgan_Lop11_Lan1
\begin{ex}%[1D9H1-3]%[ma:6a828b70]
	Hộp $A$ đựng $5$  tấm thẻ được đánh số từ $1$ đến $5$, hộp $B$ đựng $6$ tấm thẻ được đánh số từ $1$ đến $6$, hai thẻ khác nhau ở mỗi hộp đánh hai số khác nhau. Chọn ngẫu nhiên từ hộp $A$ một tấm thẻ và từ hộp $B$ hai tấm thẻ. Xét các biến cố
	\begin{itemize}
		\item $X$ là biến cố \lq\lq Chọn  được thẻ mang số lẻ từ hộp $A$\rq\rq;
		\item  $Y$ là biến cố \lq\lq Chọn được thẻ mang số chăn từ hộp $A$\rq\rq;
		\item $Z$ là biến cố \lq\lq Chọn được hai  thẻ mang số lẻ từ hộp $B$\rq\rq.
	\end{itemize}
	Tính xác suất để tích số được ghi trên ba tấm  thẻ thu được là số chẵn.
	 \shortans{0,88}
	 \loigiai{
		Ta có $X=\{1;3;5\}, Y=\{2;4\} \Rightarrow X\cap  Y=\varnothing$.\\
		Vậy $X$, $Y$  là hai biến cố xung khắc.\\
		Để có được tích ba số ghi trên ba thẻ là số lẻ thì cả ba thẻ thu được đều mang số  lẻ.\\
		Gọi $U$ là biến cố \lq\lq Tích ba số ghi trên ba thẻ  là số lẻ\rq\rq. Suy ra $\mathrm{P}\overline{U})$ là xác suất cần tìm.\\
		Dễ thấy mỗi cặp biến cố $X$ và $Z$; $Y$  và $Z$ là độc lập.\\
		 Ta có $\mathrm{P}(X)=\dfrac{3}{5}, \mathrm{P}(Y)=\dfrac{2}{5}, \mathrm{P}(Z)=\dfrac{\mathrm{C}_3^2}{\mathrm{C}_6^2}=\dfrac{1}{5}$.\\
		Khi  đó $\mathrm{P}(U)=\mathrm{P}(XZ)=\mathrm{P}(X)\cdot \mathrm{P}(Z)=\dfrac{3}{5} \cdot \dfrac{1}{5}=\dfrac{3}{25}$.\\
		 Suy ra $\mathrm{P}(\overline{U})=1-\mathrm{P}(U)=1-\dfrac{3}{25}=\dfrac{22}{25}$.
	} 
\end{ex}

% nguồn: Ok_TraLoiNgan_Lop11_Lan1
\begin{ex}%[1D9H1-3]%[ma:99924548]
	Một xạ thủ bắn lần lượt $2$ viên đạn vào một bia. Xác suất trúng đích của viên thứ nhất và viên thứ hai lần lượt là $0{,}8$ và $0{,}7$. Biết rằng kết quả  các lần bắn là độc lập với nhau. Tính xác suất để có ít nhất $1$ lần bắn trúng đích.
	 \shortans{0,94}
	 \loigiai{
		Xác suất để có ít nhất  $1$ lần bắn trúng đích là
		\[ 
		 0{,}8\times 0{,}3+0{,}2\times 0{,}7+0{,}8\times 0{,}7=0{,}94.
		\] 
	} 
\end{ex}

% nguồn: Ok_TraLoiNgan_Lop11_Lan1
\begin{ex}%[1D9H1-3]%[ma:fde2a605]
	An và  Bình không quen biết nhau và học ở hai nơi khác nhau. Xác suất để An và Bình đạt điểm giỏi về môn Toán trong kì thi cuối năm tương ứng là $0{,}92$ và $0{,}88$. Xác suất để cả An và Bình đều đạt điểm giỏi bằng bao nhiêu? (Làm tròn kết quả đến hàng phần trăm).
	 \shortans{0,81}
	\loigiai{ 
		Gọi $A$ là biến cố \lq\lq An đạt điểm giỏi về  môn Toán\rq\rq. \\
		Gọi $B$ là biến cố  \lq\lq Bình đạt điểm giỏi về môn Toán\rq\rq.\\
		Dễ thấy $A$, $B$ là hai biến  cố độc lập, khi đó $AB$ là biến cố \lq\lq Cả An và Bình đều đạt điểm giỏi môn Toán\rq\rq.\\
		Vậy xác suất để cả An  và Bình đều đạt điểm giỏi là
		 \[
		\mathrm{P}(AB)=\mathrm{P}(A)\cdot  \mathrm{P}(B)=0{,}92\cdot 0{,}88=0{,}8096 \approx 0{,}81.
		 \]
	 }
\end{ex}

% nguồn: Ok_TraLoiNgan_Lop11_Lan1
\begin{ex}%[1D9H1-3]%[ma:0c3af59f]
	Một chiếc máy có hai động cơ $I$ và $II$ hoạt động độc lập với nhau. Xác suất để động cơ $I$ chạy tốt là $0{,}8$ và xác suất để động cơ $II$ chạy  tốt là $0{,}7$. Hãy tính xác suất để cả hai động cơ đều chạy tốt là bao nhiêu?
	 \shortans{0,56}
	\loigiai{$A$ là biến cố  ``Động cơ $A$ chạy tốt''.\\
		$B$ là  biến cố ``Động cơ $B$ chạy tốt''.\\
		$C$ là biến cố ``Cả hai động  cơ đều chạy tốt''.\\
		Ta có $A$, $B$  là hai biến cố độc lập và $C=AB$.\\
		Theo công thức nhân xác suất $\mathrm{P}(C)=\mathrm{P}(A) \cdot \mathrm{P}(B)= 0{,}8 \cdot 0{,}7 =0{,}56$. 
	 }
\end{ex}

% nguồn: Ok_TraLoiNgan_Lop11_Lan1
\begin{ex}%[1D9V1-3]%[ma:2625c6a9]
	Một vận động viên thi bắn súng. Biết rằng xác suất để vận động viên bắn trúng vòng $10$ là $0{,}2$; bắn trúng vòng $9$ là $0{,}25$ và bắn trúng vòng $8$ là $0{,}3$. Nếu bắn trúng vòng $k$ thì được $k$ điểm. Vận động viên đạt huy chương vàng nếu được $20$ điểm, đạt huy chương bạc nếu được $19$ điểm và đạt huy chương đồng nếu được $18$ điểm. Vận động viên thực hiện bắn hai  lần và hai lần bắn độc lập với nhau. Xác xuất để vận động viên đạt được huy chương bạc là
	\shortans{0,1} 
	\loigiai{ 
		 Gọi $A$ là biến cố: \lq\lq Lần $1$ bắn trúng vòng $10$ và lần $2$ bắn trúng vòng $9$\rq\rq;\\
		$B$ là biến cố: \lq\lq Lần $1$ bắn trúng vòng $9$ và lần $2$ bắn trúng vòng  $10$ \rq\rq.\\
		$H$ là biến cố: \lq\lq Xạ thủ đạt  được huy chương bạc\rq\rq;\\
		Các biến cố $A$, $B$ là 2 biến cố xung khắc nên $H=A \cup B$.\\ 
		Suy ra theo quy tắc cộng ta có $P(H)=P(A)+P(B)$.\\ 
		Mặt khác $P(A)=P(B)=0{,}2\cdot 0{,}25=0{,}05$.\\ 
		Do đó $P(H)=0{,}05+0{,}05=0{,}1$. 
	} 
\end{ex}

% nguồn: Ok_TraLoiNgan_Lop11_Lan1
\begin{ex}%[1D9V1-3]%[ma:483e13d9]
	Ở ruồi giấm, tính trạng cánh dài là tính trạng trội hoàn toàn so với tính trạng cánh ngắn. Cho ruồi giấm cái cánh dài thuần chủng giao phối với ruồi giấm đực cánh ngắn thuần chủng  thu được $\mathrm{F}1$ toàn ruồi giấm cánh dài. Tiếp tục cho $\mathrm{F}1$ giao phối với nhau và thu được các con ruồi giấm $\mathrm{F}2$. Lần lượt lấy ngẫu nhiên hai con ruồi giấm $\mathrm{F}2$, tính xác suất của biến cố \lq\lq Có đúng một con ruồi giấm cánh dài trong hai con được lấy ra\rq\rq(làm tròn đến hàng phần trăm).
	\shortans{0,38} 
	\loigiai{ 
		Quy ước kiểu gene ở thế  hệ $\mathrm{F}1$ như sau
		\begin{itemize}
			\item Gene $A\colon$ ruồi  giấm cánh dài;
			\item  Gene $a\colon$ ruồi giấm cánh ngắn.
		\end{itemize}
		Ở thế hệ $\mathrm{F}2$, ba kiểu gene $AA\colon Aa\colon aa$ xuất hiện với tỉ lệ  $1\colon 2\colon 1$ nên tỉ lệ ruồi giấm cánh dài so với ruồi giấm cánh ngắn là $3: 1$.\\
		Gọi $X_1$, $X_2$ lần lượt là biến cố \lq\lq Ruồi giấm lấy ra lần thứ  nhất là ruồi giấm cánh dài\rq\rq\,và biến cố \lq\lq Ruồi giấm lấy ra lần thứ hai là ruồi giấm cánh dài\rq\rq.\\
		Ta có $X_1{,} X_2$ là hai biến cố độc lập  và $\mathrm{P}(X_1)=\mathrm{P}(X_2)=\dfrac{3}{4}$.\\
		Xác  suất của biến cố \lq\lq Có đúng một con ruồi giấm cánh dài trong hai con được lấy ra\rq\rq\,là
		\begin{eqnarray*}
			\mathrm{P}(X_1\overline{X_2}\cup\overline{X_1}X_2) 
			&=& \mathrm{P}(X_1\overline{X_2})+\mathrm{P}(\overline{X_1}X_2)\\ 
			&=& \mathrm{P}(X_1) \mathrm{P}(\overline{X_2})+P(\overline{X_1})  P(X_2)\\
			 &=& 2\cdot \dfrac{3}{4}\cdot \dfrac{1}{4}\\
			&=&  \dfrac{3}{8}.
		\end{eqnarray*}
	} 
\end{ex}

% nguồn: Ok_TraLoiNgan_Lop11_Lan1
\begin{ex}%[1D9V1-3]%[ma:e0c7a1a2]
	Một bình đựng $5$ viên bi xanh và $3$ viên bi đỏ (các viên bi chỉ khác nhau về màu sắc). Lấy ngẫu nhiên một viên bi, rồi lấy ngẫu nhiên một viên bi nữa. Tính xác suất của biến cố \lq\lq Lấy lần thứ hai được một viên bi xanh\rq\rq\, (làm tròn  kết quả đến hàng phần trăm).
	\shortans{0,63} 
	 \loigiai{
		Gọi $A$ là biến cố \lq\lq Lấy lần thứ hai được một viên bi xanh\rq\rq .\\ 
		Gọi $B$ là biến cố \lq\lq Lấy lần thứ  nhất được bi xanh, lấy lần thứ hai cũng được một bi xanh\rq\rq .\\
		Xác suất biến cố $B$ là  $\mathrm{P}(B)=\dfrac{5}{8}\cdot\dfrac{4}{7}=\dfrac{5}{14}$.\\
		 Gọi $C$ là biến cố \lq\lq Lấy lần thứ nhất được bi đỏ, lấy lần thứ hai được bi xanh\rq\rq .\\
		Xác suất biến  cố $C$ là $\mathrm{P}(C)=\dfrac{3}{8}\cdot\dfrac{5}{7}=\dfrac{15}{56}$.\\
		Ta có $A=B\cup C$, vì $B$, $C$ là hai biến cố xung khắc nên ta có\\ 
	$\mathrm{P}(A)=\mathrm{P}(B\cup  C)=\mathrm{P}(B)+\mathrm{P}(C)=\dfrac{5}{14}+\dfrac{15}{56}=\dfrac{35}{56}=\dfrac{5}{8}$.}
\end{ex}

% nguồn: Ok_TraLoiNgan_Lop11_Lan1
\begin{ex}%[1D9V1-3]%[ma:1072d037]
	Gieo hai con súc sắc I và II cân đối, đồng chất một cách độc lập. Ta có biến cố $A$: \lq\lq Có ít nhất một con súc sắc xuất hiện mặt $6$ chấm\rq\rq . Lúc  này giá trị của $\mathrm{P}(A)$ là
	(kết quả làm tròn đến hàng  phần trăm).
	 \shortans{0,31}
	\loigiai{ 
		Gọi $A_i(i=1;2)$ là biến cố: \lq\lq Con súc sắc  thứ $i$ ra mặt $6$ chấm\rq\rq \\
		$ \Rightarrow A_1 $ và $A_2$ là hai biến cố độc lập và  ta có $\heva{&\mathrm{P}(A_1)=\dfrac{1}{6}\\&\mathrm{P}(A_2)=\dfrac{1}{6}.}$ \\
		 Thay vì tính $\mathrm{P}(A)$ ta đi tính $\mathrm{P}\left(\overline{A}\right)$. Ta có $\overline{A}=\overline{A_1}\cdot\overline{A_2}$.\\
		$\mathrm{P}\left(\overline{A}\right)=\mathrm{P}\left(\overline{A_1}\right)\cdot\mathrm{P}\left(\overline{A_2}\right)=\left(1-\mathrm{P}(A_1)\right)\cdot\left(1-\mathrm{P}(A_2)\right)=\dfrac{5}{6}\cdot\dfrac{5}{6}=\dfrac{25}{36}$.\\ 
	 Vậy $\mathrm{P}(A)=1-\mathrm{P}\left(\overline{A}\right)=1-\dfrac{25}{36}=\dfrac{11}{36}$.}
\end{ex}

% nguồn: Ok_TraLoiNgan_Lop11_Lan1
\begin{ex}%[1D9V1-3]%[ma:483c2865]
	Một bệnh truyền nhiễm có xác suất lây bệnh là $0{,}8$ nếu tiếp xúc với người bệnh mà không đeo khẩu trang; là $0{,}1$ nếu tiếp xúc với người bệnh mà có đeo khẩu trang. Chị Hoa có tiếp xúc với người bệnh hai lần, một lần đeo khẩu trang  và một lần không đeo khẩu trang. Tính xác suất để chị Hoa bị lây bệnh từ người bệnh truyền nhiễm đó.
	\shortans{0,82} 
	\loigiai{ 
		Gọi $A$ là biến cố: ``Chị Hoa bị nhiễm bệnh khi tiếp xúc người bệnh mà không đeo khẩu trang'' và $B$: ``Chị Hoa bị nhiễm bệnh khi tiếp  xúc với người bệnh dù có đeo khẩu trang''.\\
		Dễ thấy $\overline{A}, \overline{B}$ là hai biến cố độc lập.\\ 
		Xác suất để chị Hoa không nhiễm bệnh trong cả hai lần tiếp xúc với người bệnh là $\mathrm{P}(\overline{A}\cap\overline{B})=\mathrm{P}(\overline{A})  \cdot \mathrm{P}(\overline{B})=0{,}2 \cdot 0{,}9=0{,}18$.\\
		Gọi $\mathrm{P}$ là xác suất để chị Hoa  bị lây bệnh khi tiếp xúc người bệnh, ta có
		\[\mathrm{P}=1-P(\overline{A}\cap\overline{B})=1-0{,}18=0{,}82.\] 
	} 
\end{ex}
